% 本文件由 LaTeX 排版 Agent 根据有效 translation.json 自动生成。
% author=Augustine of Hippo; work_id=1303; title=论教导未受教育者

\chapter{论教导未受教育者}

% structure: p0001 source=h1 target=omitted reason=page-title-already-rendered-by-parent

% structure: p0002 source=h2 target=section reason=relative-heading-rank; base=h2

\section{序言}

在\WorkTitle{Retractations}第二卷第十四章中，\Person{Augustine}作了如下陈述：\Quote{我们还有一本论\WorkTitle{教导未受教育者}的著作，[即《教导未受教育者》，或称\WorkTitle{De Catechizandis Rudibus}]，这确实是它的明确标题。在这本书中，当我说：\InnerQuote{那位天使，连同一群作为他追随者的其他神灵，骄傲地背弃了天主的服从，成了魔鬼，对天主并无伤害，而是伤害了自己；因为天主知道如何处置离开祂的灵魂}，更恰当的说法应是：\InnerQuote{离开祂的神灵}，因为所讨论的是天使。这本书的开头是这样的：\InnerQuote{弟铎格拉恰斯兄，你曾请求我。}}\par

这篇在引用段落中描述的著作，是\Person{Augustine}在公元400年左右与其他手头著作一起审阅的，因此可认为属于那个时期。据推测，收信人可能与司铎\Person{Deogratias}是同一人，据我们读到的现在列为第一百零二封的书信，\Person{Augustine}大约在406年写信给\Person{Deogratias}，答复从\Place{Carthage}转寄给他的异教徒的一些问题。\par

本笃会编辑们对该论著的介绍如下：\Quote{应\Place{Carthage}一位执事的请求，\Person{Augustine}承担了教导教理讲授技巧的任务；首先，他给出了一些指示，指出这种职责不仅要有固定的方法和适当的顺序，而且要避免冗长乏味，并怀着愉快的心情来履行。随后，他将这些指示付诸实践，通过发表两篇互相平行的演讲来示范他的意思，一篇稍长，一篇非常简短，但都适合用来教导任何渴望成为基督徒的人。}\par

% structure: p0006 source=h2 target=section reason=relative-heading-rank; base=h2

\section{第一章——\Person{Augustine}如何应\Place{Carthage}一位执事的请求而写作}

1. \Person{Deogratias}兄，你曾请求我，要我寄给你一些书面的东西，可以在教导未受教育者方面对你有所帮助。因为你告诉我，在\Place{Carthage}，你担任执事的地方，由于你享有\OriginalTerm{教理讲授}{catechising}方面的丰富恩赐——既源于对信仰的深入了解，又源于引人入胜的论述方式——因此经常有人被带到你面前，需要从最基础开始学习基督信仰；但是你几乎总是遇到困难，不知如何适当地阐述那构成我们基督徒的信仰的确切教义：关于我们的陈述应从何处开始，应延续到何处；以及在叙述结束后，是否应当使用某种劝勉，还是仅仅列出那些听我们讲话的人所应遵守的诫命，借此他才能知道如何保持基督徒的生活和身份。同时，你坦白并抱怨说，你经常遇到这样的情况：在冗长而乏味的讲话中，你不仅对学习者——你试图通过言语教导他——以及对在场的其他听众毫无益处，甚至令自己也感到厌恶；正是这些困境迫使你将这份我爱你的重担加于我身上，使我在百忙之中不觉得写些关于这个主题的东西是一件麻烦事。\par

2. 至于我自己，如果在我通过我主的恩赐所能呈现的能力中，同样我主要求我向那些祂已使我成为弟兄的人提供任何帮助，我感到不仅出于我对你所负的基于亲密友谊的爱与服侍，而且出于我普遍欠我母亲教会的债，我绝不能拒绝这项任务，反而要以迅速而虔诚的意愿承担起来。因为我越希望主的宝藏被广泛分配，就越有责任，如果我发现那些作为我同仆的管家在分配时遇到困难，就尽我所能，使他们能够轻松而迅速地完成他们勤勉而热切追求的目标。\par

% structure: p0009 source=h2 target=section reason=relative-heading-rank; base=h2

\section{第二章——论为何常发生听者愉悦而讲者反感的事，并对这一事实作何解释}

3. 然而，关于你在这种努力中私下抱持的想法，我不愿你因自己常觉得所发表的言论贫乏而乏味而感到困扰。很可能在你想教导的人看来，情况并非如此，而只是因为你热切渴望有更好的内容可听，才觉得自己的话不配入他人之耳。事实上，我也几乎总是对自己的演讲感到不满。因为我贪求更好的东西，在我开始用可理解的言语将其表达出来之前，我常常在内心里享受它的拥有；然而，当我表达的能力证明不及我内心的领悟时，我便因舌头未能回应心灵而悲伤。因为我希望听我的人能像我一样对那题目有完全的理解；但我觉察到我的讲话方式未能产生那种效果，这主要源于如下情况：理智的领悟像一道闪电般迅速在心灵中扩散，而言语则缓慢、耗时，且性质截然不同，因此当言语移动时，理智的领悟已退回其隐秘的居所。然而，由于它以奇妙的方式在记忆中留下了某些自己的印记，这些印记随着音节的短暂停顿而持续；并且凭借这些印记，我们形成了可理解的符号，这些符号被称为某种语言，例如拉丁语、希腊语、希伯来语或其他语言。这些符号可以是思想的对象，也可以实际由声音发出。另一方面，这些印记本身既非拉丁语、希腊语、希伯来语，也非任何其他民族所特有，而是在心灵中像容貌在身体上一样成为现实。因为愤怒在拉丁语中用一个词表示，在希腊语中用另一个词，在其他语言中用不同的词语，各有不同。但愤怒之人的面容既非（特指）拉丁语，也非（特指）希腊语。因此，当某人说\Emph{Iratus sum}，并非每个民族都明白，而只有拉丁人明白；然而，如果他的心灵在激起怒气时的情绪展现在脸上并影响容貌，那么所有看见他的人都会明白他在生气。但是，我们无法将理智领悟在记忆上留下的那些印记取出来，并像容貌那样清晰而明显的形式，通过声音来展现给听者感知。因为前者在心灵之内，而后者在身体之外。因此我们必须猜想，我们口中的声音与理智的冲击相差多远，因为它甚至与记忆产生的印象也不相符。现在，我们常有一种情况：在渴望听众获益时，我们试图像内心的理智领悟那样向他表达自己，而正在努力时我们却缺乏说话的能力；然后，因为未能成功，我们感到懊恼，并在疲倦中憔悴，仿佛在从事徒劳的劳动；而正因这疲倦，我们的言论变得比最初引发这种乏味感时更加无力而无意义。\par

但那些渴望听我说话的人的认真态度时常向我显示，我的言辞并不像我自己感觉的那样冷淡。从他们表现出的喜悦中，我也看出他们从中获得了某种益处。我勤勉地致力于不让自己在为他们提供一项我所察觉他们欣然接受的服务时有所缺失。同样，从你这边来看，那些需要接受信仰教导的人频繁地被带到你面前，这一事实本身应当帮助你理解，你的言论并非像你自己感到的那样令他人不悦；你也不应仅仅因为你未能以自己所期望的方式陈述你明辨的事物而认为自己不结果实；因为，或许你也同样无法按照你所希望的方式去明辨它们。因为在此生，除了\BibleQuote{如同谜语和透过镜子}一样，谁又能看见呢？爱本身也没有足够的力量撕裂肉身的黑暗，渗入那永恒的宁静——连转瞬即逝的事物也从那里获得它们所闪耀的光辉。然而，既然善人日复一日地向那超越旋转穹苍、未经黑夜侵入的视景迈进——\BibleQuote{眼所未见，耳所未闻，人心也未曾想到的}——当我们从事教导未受训者时，导致我们的言论在自己估量中变得毫无价值的理由，没有比这更大的了：即我们以非凡的方式明辨是一种乐趣，而以普通的方式言说却是件苦差。事实上，当我们自己也乐在其中时，听众会感到更大的满足；因为我们演讲的脉络会受到我们所感受的欢愉的影响，从而更顺畅、更受欢迎地从我们口中流出。因此，对于那些被推荐为信仰条文的事项，并非难事去颁布指示：关于叙述应从哪里开始和结束，或关于叙述应如何变化，以便有时更简短、有时更详尽，却始终充实完善；以及关于哪些场合适合使用简短形式，哪些场合宜采用较长形式。但关于如何做到这一切，使每个人在讲授教理时都能乐于其工作（因为他做得越好，就越有吸引力）——这需要最大的考量。然而，我们不需要远求即可获得在这一领域起作用的准则。因为，如果在属肉体的事物上，\BibleQuote{天主喜爱乐意捐献的人}，那么在属神的事物上岂不更加如此？但我们确信这喜乐能在适当时候与我们同在，则依赖于赐予我们这些规诫的那位的仁慈。因此，按照我对你愿望的理解，我们将首先讨论叙述方法，然后讨论颁布指示和劝勉的职责，之后讨论如何获得前述的喜乐，尽天主所赐予我们的思想。\par

% structure: p0012 source=h2 target=section reason=relative-heading-rank; base=h2

\section{第三章 论教导未受训者时所应采用的完整叙述}

叙述是完整的，当每个人在最初受训时，从经文所写的\BibleQuote{在起初天主创造了天地}开始，一直讲到教会的现时代。但这并不意味着，我们应当凭记忆逐字重复整个梅瑟五书、整个民长纪、列王纪和厄斯德拉纪、整个福音和宗徒大事录，即使我们已经全部记住；也不意味着我们应当将所有这些卷册中所有内容用自己的话复述，并以那种方式整体展开和解释。因为时间不允许，也没有任何必要强求如此。我们所应当做的是，对一切事项作综合扼要的概述，以便可以选出某些更奇妙的事实——这些事实以超乎寻常的满足感被聆听，并且被列为历史的精确转折点；这样，我们不应（如果可以这么说的话）只让它们在包裹中显现、随即从我们视线中夺走，而应当在它们上面停留一段时间，从而仿佛展开它们、打开它们以供观看，并将它们作为需要审视和赞叹的事物呈现在听众心中。至于所有其他细节，则应快速略过，并以此种方式引入并编织入叙述中。采用这一计划的效果是，我们特别希望引起注意的那些特定事实，因其他事实让道而获得了更突出的地位；同时，我们急切希望以论述激发兴趣的初学者，不会带着已经疲惫的头脑接触这些主题，也不会给本应通过我们教导而受教的记忆带来混乱。\par

6. 实际上，在一切事上，我们不仅当将自己的目光固定在诫命的终极，即\BibleQuote{出于纯洁的心、良好的良心和无伪的信德的爱德}，而我们所说的一切都当指向它；同样，我们以言语教导之人的目光也当归向同一目标，并被引导至那方向。事实上，在主的来临之前，圣经中所记载的一切，其写作原因无他，正是为了这一点——即祂的来临应被铭记，而将要形成的教会，即万民中天主的子民，应被预先宣告；这教会就是祂的身体，凡在祂来临之前活于这世界，并相信祂将来临的一切圣人，也都与此身体联合并被纳入，正如我们相信祂已来临一样。（打个比方）雅各伯在诞生时，先从母腹伸出一只手，这只手还抓住了在出生中领先于他的兄弟的脚；接着头出来了，然后最后，自然而然，其余肢体也出来了；然而，头在尊贵和能力上不仅超越其后出来的肢体，甚至超越在出生过程中先于它的手，它真正是首，虽然不是显现的时间上，至少在本性的次序上是如此。同样，主耶稣基督，在祂显现于肉身之前，并以某种方式从祂隐密的母腹中出来，在世人面前以人的形象出现，天主与人之间的中保，\BibleQuote{祂是在万有之上，永远受赞美的天主}，在圣祖和先知们身上，预先把祂身体的一部分——正如通过一只手——预先发出了祂自己即将诞生的记号，并且用法律之束缚（如同祂的五指）压制了那些在骄傲中先于祂的百姓。因为通过五个时代，对祂注定来临的预告和预言从未间断；与此一致的是，藉着祂赐下法律的那一位写了五部书；而骄傲的人，属血肉的，企图\BibleQuote{建立自己的义}，并未因基督张开的手而得祝福，反而因那压缩闭合的手被剥夺了这种福分；因此，他们的脚被绑住，\BibleQuote{他们已颠覆摔倒，我们却已复活，挺立站住}。但是，尽管正如我所说，主基督确实在那些在出生时间上先于祂的圣人身上预先发出了祂身体的一部分，然而祂自己是身体——教会——的头，而所有这些人都因相信他们所预言宣告的祂，而被连接于那同一个以祂为头的身体。因为他们没有因在祂之前完成自己的路程而与那身体分离，反而因他们的服从而与之合而为一。因为手虽可能在头之前伸出，但它仍与头之下方相连。因此，凡先前所写的，都是为教训我们而写的，都是我们的预像，且发生在这些人身上作为预像。而且，它们是为我们这些万世的结局已经来到的人而写的。\par

% structure: p0015 source=h2 target=section reason=relative-heading-rank; base=h2

\section{第四章　基督来临的伟大原因乃是爱的显扬}

7. 再者，还有什么比这更明显的理由来说明主的来临，就是天主在我们内显示祂的爱，强力地显示它，因为\BibleQuote{当我们还是罪人的时候，基督为我们死了}？而且，这正意图在于，既然爱德是\BibleQuote{诫命的终极目标}和\BibleQuote{法律的满全}，我们也能彼此相爱，并为弟兄舍命，正如祂为我们舍命一样。而且关于天主本身，其目的在于，即使爱祂曾经是一件令人厌烦的任务，如今至少我们对祂回报以爱不再会感到厌烦，因为\BibleQuote{祂先爱了我们}，并且\BibleQuote{没有怜惜自己的独生子，反而为我们众人交付了祂}。因为再没有比预先表现出爱更有力的爱的邀请了；那灵魂真是过于坚硬，既然不愿意付出爱，也不愿意回报爱。但是，即使在邪恶而污秽的爱情中，我们看到那些渴望被爱回报的人如何将其作为特别而全神贯注的事务，通过他们能力范围内的种种证明，使他们自己爱情的力量明白可见；我们也看到他们如何装作在他们所提供的东西中表现出一种正义的外表，这种外表或许能使他们以某种方式要求那些他们努力诱惑的灵魂公平地回应他们；此外，当他们观察到他们急于占有的心也被同样火焰所激动时，他们的热情燃烧得更炽烈：如果这样，我再说一遍，那先前麻木的灵魂一旦感到自己被爱就立刻被激发，而那已经燃起的灵魂一旦得知自己的爱情得到了回报就变得更加火热，那么显而易见，没有什么比这种情境更能说明爱情为何被激发或扩大了：即一个尚未爱过的人得知自己成为了爱的对象，或者一个已经爱着的人要么希望被爱回报，要么已经有了爱情得到回应的证据。如果这在卑劣的爱情中尚且成立，那么在真友谊中岂不更是如此？因为在这涉及伤害友谊的问题中，我们所应当谨慎注意的，难道不是不让我们的朋友有任何理由认为我们根本不爱他，或认为我们爱他不如他爱我们吗？如果，他确实被引导到这种信念，他就会对那种人们彼此享受亲密交流的爱情冷淡下来；如果他没有那种软弱的性格，即这种对感情的伤害会完全冻结他的爱，那么他仍然将自己局限于那种爱的方式：他爱，不是为了自己的享受，而是为了成全他人的益处。但是，也值得我们留意：尽管上级也希望被下级爱，并且因他们殷勤的关注而满意，而且他们越是知道这一点，就越对这些下级怀有更大的爱；然而，下级一旦知道自己被上级所爱，又是怀着何等强大的爱燃烧起来！因为我们在那里看到更为感恩的爱情，它并不在需求的干涸中消耗自己，而是在慷慨的丰沛中流淌。因为前一种爱情是悲惨的，后一种则是怜悯的。而且，如果下级甚至对自己可能被上级所爱感到绝望，那么当上级出于自己的意志屈尊显示他多么爱这个人时，这个人将无法表达地被感动而爱，而这个人原本绝不敢奢望如此伟大的善。但是，有谁在审判者的身份上比天主更优越？有谁在罪人的身份上比人更绝望？——我问，这个人曾经更加毫无保留地将自己交给那些骄傲的权势的监护和统治，这些权势不能使他幸福，因为他更加绝对地绝望于自己可能成为那个不愿意在邪恶中被高举、而是在善中被高举的权势所感兴趣的对象。\par

8. 因此，如果基督降临的主要目的就是让人明白天主多么爱他，并让他明白这一点，为的是使他被那首先爱他的爱所点燃，并因那成为我们近人的主的命令和榜样而爱邻人——主成为我们的近人，乃是因他爱了人，而人原本不是他的近人，远远地漂泊在外。如果全部古代所写的圣经，都是为了预指主的降临而写的；如果此后凡凭天主权威写下的记载，都是基督的记录，并提醒我们关于爱，那么显而易见，爱天主和爱邻人这两条诫命，不仅是一切法律和先知的基础——当时主说这话时，法律和先知还是唯一的圣经——而且也是后来为了我们的健康而写就并流传下来的全部神圣文学的基础。因此，在旧约中隐藏着新约，在新约中彰显着旧约。根据那隐藏，属血肉的人以血肉的方式理解，无论古今，都处于惩罚性的恐惧之下。反之，根据这彰显，属神的人——我们视之为包括当时以虔诚叩门而得以发现隐藏之事的人，以及现在不以骄傲之心寻求、免得敞开的向他们也关闭的人——以属神的方式理解，藉着所赐予的爱而获得了自由。因此，既然没有什么比嫉妒更与爱相悖，而骄傲是嫉妒之母，那么主耶稣基督，既是天主又是人，就既是天主对我们的爱的彰显，也是与我们同在的人性谦卑的榜样，为要使我们的巨大肿胀因更大的反治疗法而痊愈。因为人骄傲真是极大的不幸！但天主谦卑却是更大的仁慈！因此，请将爱作为你面前的目标，你一切言谈都要指向它，无论你叙述什么，都要以这样的方式叙述，使你对他讲话的人听了就相信，信了就盼望，盼了就爱。\par

% structure: p0018 source=h2 target=section reason=relative-heading-rank; base=h2

\section{第五章 前来接受教义讲授者，应在渴望成为基督徒时，就其观点受到考查}

9. 此外，正是在天主的这种严厉中——它以一种最有益的恐惧震动必死之人的心——爱得以建立起来；如此，人因受到他所畏惧的那位的爱而欢喜，便勇敢地去回报爱，同时却又害怕得罪那爱他的主，即使他能够不受惩罚。因为确实很少有，不，我宁可说从未有人，在没有被某种天主之恐惧击中之前，就来向我们表示愿意成为基督徒的。因为，如果一个人是因为期待从人那里得到某种好处（他以为自己别无他法取悦于人），或是为了逃避人的不满或敌意（他极为害怕）而想成为基督徒，那么他想要成为基督徒的愿望并不如他想要假装基督徒的愿望那样真诚。因为信仰不在于服从的身体，而在于相信的心灵。但毫无疑问，常常是天主的慈爱通过教义讲授者的服务临到，使那人受讲话的影响，现在真心渴望成为他本来决心只是假装的那样。一旦他开始有这种渴望，我们就可以判断他真正来到了我们面前。的确，我们感知到某人身体已经在场，而他的心真正来到我们面前的确切时刻，对我们是隐藏的。尽管如此，我们还是应当这样待他，使这种愿望在他心中产生，即使此刻它尚未存在。因为这样的劳作不会白费，因为，即使有丝毫愿望，我们的行动也肯定会加强它，尽管我们不知道它开始的时间或时辰。如果可能，从了解那人的人那里预先了解他的心态，或促使他前来接受信仰的原因，这当然是有用的。但如果没有其他人可以提供这些信息，那么确实应当询问那人本人，以便从他回答的话中引出我们讲话的开端。如果他是怀着虚伪的心来的，只贪求人的好处或想逃避不利，那么他肯定会说假话。然而，我们应当从他所说的假话出发。但这不应带着公开揭穿他谎言的意图去做，好像你已经认定是假话；而是假定他声称自己怀有一种真正值得称赞的目的（无论这声称是真是假），我们应当称赞和表扬他回答时所宣称的那种目的，以便使他感到：实际成为他想要看起来的那种人是一种快乐。另一方面，假设他给出的观点不同于应当存在于一个将受教于基督教信仰的人心中的观点，那么就要以异乎寻常的仁慈和温和责备他，视他为无知的初学者，简明而严肃地指出并赞扬基督教教义的真实目的，而且以这样的方式行事：既不预先叙述（那应在之后给出），也不冒然将这类陈述强加于尚未准备好的心灵，这样你就可以引导他渴望那到目前为止，他出于错误或伪装而尚未渴望的事。\par

% structure: p0020 source=h2 target=section reason=relative-heading-rank; base=h2

\section{第六章　——　论如何开始要理讲授，以及从世界受造之日起直至教会现今时代的历史事实的叙述}

10. 但如果他的回答是，他曾遇到某种神圣的警告或某种神圣的恐惧，促使他成为基督徒，那么这就为我们的讲论提供了一个最令人满意的开端，因为它暗示了天主对我们多么关怀。然而，他的思想当然应当从这类事物——无论是奇迹还是梦——转向更坚实的道路和圣经更可靠的圣言；这样他也可以明白，在他投身于圣经之前，那警告是如何仁慈地预先给予他的。确实应当向他指出，主自己既不会这样劝告他、催促他成为基督徒并加入教会，也不会用这样的迹象或启示教导他，除非祂愿意为了他更大的安全和保障，让他走上一条已在圣经中预备好的道路，在这条路上他不应寻求看得见的奇迹，而应学会盼望不可见之事的习惯，并且也应在清醒时而非睡梦中接受警告。此时应当开始叙述，从天主创造了所有美好的事物这一事实开始，并如我们所说，一直延续到教会现今的时代。这应当以这样的方式进行：为我们所叙述的每一件事和每一个事件，给出其原因和理由，使我们可以将它们分别指向那爱德的终向，无论是正在做某事的人的眼目，还是正在说话的人的眼目，都不应偏离那终向。因为，即使在处理诗人的神话——这些不过是虚构的创造和为了以琐事为食粮的心灵的愉悦而编造的东西——时，那些享有良好声誉和名声的语法家仍努力将它们引向某种（假定的）用处，尽管这用处本身可能只是空虚的、严重偏向于这世界的粗俗营养：那么，我们岂不是更应当谨慎，以免我们叙述的真正真理，由于对其原因缺乏深思熟虑的说明，而被接受为一种无益于实际好处的满足，或者更糟，被接受为一种可能有害的贪婪？同时，我们也不应以这样的方式陈述这些原因，以至于离开我们叙述的正当进程，让我们的心灵和舌头陷入更复杂讨论的棘手问题中。但我们所提出的解释的简单真理，应当像连接一串宝石的金线一样，既不影响装饰的优美对称，也不因自身不当的插入而干扰它。\par

% structure: p0022 source=h2 target=section reason=relative-heading-rank; base=h2

\section{第七章　——　论在此叙述之后应阐述的复活、审判及其他主题}

在完成这番叙述之后，应阐述复活的希望，并视听者的能力和力量以及我们可支配的时间，论证反对那些不信者对身体复活的虚妄嘲弄，以及对未来审判的辩论——这审判对善者彰显其良善，对恶者显示其严厉，对所有人表明其公义。在痛恨和惊惧地宣告不敬者的惩罚之后，接着应谈论义人忠信者的王国、那超乎诸天的城邑及其喜乐。此外，我们还应装备并激励人弱小的力量，以抵挡诱惑和绊脚石，无论这些是来自教会外部，如反对外邦人、犹太人、异端者，还是来自教会内部，如反对\OriginalTerm{主的打谷场上的糠秕}{chaff of the Lord's threshing-floor}。但这并不意味着我们要与每种悖逆之人逐一辩论，并用一套套论证驳斥他们所有错误的观点；相反，我们应在有限的时间内，以合宜的方式指出这一切都已预先宣告，并说明诱惑在训练信徒中有何益处，以及天主的忍耐——他决定容忍它们直到终结——的榜样带来何等安慰。此外，当听者被装备以对抗这些对手（他们的悖逆群体虽在身体上充满教会）时，也应同时简明适当地阐述基督徒和可敬生活方式的规诫，以免他轻易被那些酗酒、贪婪、欺诈、赌博、通奸、行淫、喜爱公共表演、佩戴不洁符咒、行巫术、占星术或任何虚妄邪恶之术的占卜者，以及其他类似之人引入歧途；也不要因看见许多被称为基督徒的人喜爱这些事、参与其中、为之辩护、推荐它们，甚至说服他人使用它们，而自以为可以肆无忌惮地效法。因为对于持续这种生活方式的人所注定的结局，以及他们在教会内（最终将被分离出去）应当如何被容忍——这些话题应用圣经的见证来教导学习者。然而，他也应事先得知，教会中有许多善良的基督徒，他们是天上耶路撒冷最真实的公民——如果他自己也愿意成为这样的人。最后，必须恳切地警告他，不要把希望寄托在人身上。因为人无法轻易判断谁是义人。即使这件事容易做到，但义人的榜样摆在我们面前，其目的也不是要我们藉他们成义，而是藉效法他们，明白我们自己如何被那使他们成义者成义。这事的结果必将获得最高赞许——即当听我们（其实是听天主藉着我们说话）的人开始在道德和知识上有所进步，并以热忱踏上基督的道路时，他不会将这改变归功于我们或自己，而是爱他自己、爱我们，以及他所爱的一切朋友，都在他内、并为了他的缘故——那曾爱了他这位敌人，为使他成义并成为朋友的那一位。既然我们已经讲到这里，我认为你不需要任何导师告诉你，当你的时间或听者时间有限时如何简要讨论，以及当你有更多时间时如何更详尽地论述。因为事情本身的需要自然会提示这一点，无需任何建议者的劝告。\par

% structure: p0024 source=h2 target=section reason=relative-heading-rank; base=h2

\section{第八章 对受过通识教育者进行要理讲授的方法}

12. 然而，另有一种情形显然不容忽视。我是指有人前来接受\OriginalTerm{要理讲授}{catchetical instruction}，他曾在\OriginalTerm{博雅之学}{liberal studies}方面受过栽培，早已决心成为基督徒，并专为成为基督徒而投奔你。几乎必然会有这样的情形：这样一个人已经对我们的圣经和文献有了相当的认识；带着这些知识，他前来可能只是为了领受圣事。因为这类人的习惯是在成为基督徒之前（而非当时）就仔细探究一切事情，并且早早地与任何他们能接触到的人交流、讨论自己内心的感受。因此，对他们应采取一种精简的方法，以免令人厌烦地灌输他们已知的事情，而应轻轻带过。因此我们应说，我们相信他们已经熟悉这样或那样的事，因而我们只是简单扼要地列举那些需要正式向未受教导和无知的人强调的事实。我们应努力如此进行，以使这位有学识的人若先前已了解我们的任一主题，现在不会像听老师讲授那样听到它；假若他对其中任何事尚不知情，他仍能在我们回顾他所熟悉的事情时学习到这些。此外，询问他本人是出于什么动机渴望成为基督徒，无疑是有益的；这样，你若发现他是因书籍（无论是正典著作还是值得研读的文人作品）而作出这个决定，就可以在开始时对这些书说些什么，按照它们各自在正典权威和解经者巧妙勤勉方面所具有的不同长处，适当地表示赞赏；对于正典圣经，尤其要赞美其中在惊人的崇高主题旁所显示的极有益处的谦逊（措辞）；而在其他作品中，根据每位作者特有的才能，则要注意一种更响亮、更圆润的雄辩风格，这种风格适合那些更骄傲因而更软弱的心灵。我们当然还应引导他谈谈他自己，以便让我们了解他主要读了哪位作者，更熟悉哪些书，因为这些书促使他渴望加入教会。当他告诉我们这些信息后，若上述书籍为我们所知，或者我们至少有教会的报告作为依据，认为它们是由某位著名的公教学者所写，我们应高兴地表示赞许。但若相反，他偶然读到某些异端者的著作，并可能在无知中保留了真信仰所谴责的任何东西，而仍以为那是公教教义，那么我们必须勤勉地教导他，向他展示（以其正当的优越性）普世教会的权威，以及许多其他最博学的、以其论辩和著作维护真理而闻名之人的权威。同时，必须承认，即使是那些作为真诚信徒离世并留给后世一些基督徒著作的人，在他们小作品中的某些段落，或因未被理解，或因（如人性之常）心智眼目无力深入更深奥的奥秘，在追求真理表象时偏离了真理本身，从而为傲慢放肆之人构建和产生某些异端提供了机会。这并不奇怪，因为即使在正典著作本身（其中一切表达都最为稳妥）中，我们也看到，有些人并非仅仅因为某些段落的理解与作者原意或真理本身不一致（若仅如此，谁不愿欣然宽恕那愿意接受纠正的人性软弱呢？），而是因为执拗地以最激烈的愤慨和无耻的傲慢为自己所持的乖谬错误之见解辩护，从而导致了众多有害教义的产生，撕裂了基督徒共融的合一。所有这些主题，我们都应在和缓的交谈中与那位走近基督徒团体的人讨论——他不是以所谓未受教育者的身份，而是以在博学书籍中受过完善文化训练者的身份。在嘱咐他提防骄傲的错误时，我们应仅采取那使他前来投奔我们的谦卑所容许的权威。此外，按照救恩教义的准则，其他需要叙述或讨论的事，无论是涉及信仰、道德生活的问题，还是关于诱惑的问题，都应按我指示的方式逐一处理，并应参照那更卓越的道路（已提及）。\par

% structure: p0026 source=h2 target=section reason=relative-heading-rank; base=h2

\section{第九章　论当如何对待文法家与职业演说家}

13. 还有一些来自最低级的文法学校和职业演说家的学生，你不敢将他们列为未受教育者，也不属于那些在重大问题中锻炼心智的非常博学的阶层。因此，当这样的人——他们在言谈艺术上似乎优于其余人类——怀着成为基督徒的意向接近你时，在你与他们交流时，你要比对待那些其他未受过教育的听众更加尽力地表明：应当勤奋地劝诫他们穿戴基督徒的谦逊，并学会不轻视那些他们发现更谨慎地避免行为上的恶习而非言语上的过失的人；而且，他们不应过分自信，以至于将训练有素的舌头与纯洁的心灵相提并论，而他们往往习惯将舌头置于优先地位。但最重要的是，应当教导这样的人聆听圣经，使他们既不会仅仅因为朴实的雄辩没有浮夸就轻视它，也不会认为我们在那些书中读到的人的言语和行为，被包裹在\OriginalTerm{肉体的包裹}{carnal wrappings}中，就不应当为了（充分）理解它们而揭开并阐明，而仅仅按照字句的声音来理解。至于同一隐藏意义的有用性——正因为这种意义的存在，它们才被称为奥秘——这些\OriginalTerm{谜语式言论}{enigmatical utterances}的错综复杂在激发我们对真理的热爱和驱散厌倦的麻木方面所具有的力量，正是这些人必须通过实际经验来体会的：当某个题目赤裸裸地摆在他们面前时未能打动他们，而通过详细阐发\OriginalTerm{寓意意义}{allegorical sense}，其意义被引发出来。因为让这样的人知道思想应当优先于言辞，正如灵魂优先于身体，这是极为有益的。由此也得出，他们应当渴望聆听以真理著称的言论，而非以雄辩著称的言论；正如他们应当渴望拥有以智慧著称的朋友，而非以美貌著称的朋友。他们还应明白，除了灵魂的情感，没有声音能进入天主的耳朵。因为这样，如果他们碰巧观察到教会中的某些主教或圣职人员以充满\OriginalTerm{野蛮用语和语法错误}{barbarisms and solecisms}的言语呼求天主，或者未能正确理解自己所念的词语并混淆停顿，他们就不会嘲笑。当然，这并不意味着不应纠正这些错误，以便人们能向他们所清楚理解的内容说\Quote{Amen}；而是说，那些已经明白的人应当虔诚地容忍这些事，正如在法庭上，言语的恩宠在于声音，而在教会中，言语的恩宠在于情感。因此，法庭的言语有时可称为好的言语，但绝不是恩宠的言语。此外，关于他们将要领受的圣事，对于较聪明的人，只需听其意指什么就够了。但对于智力较迟钝者，则需要采用较为详细的解释，并加上比喻，以免他们轻视所见之物。\par

% structure: p0028 source=h2 target=section reason=relative-heading-rank; base=h2

\section{第十章 论在教理讲授中获得喜乐，以及导致慕道者厌倦的各种原因}

14. 在此你或许希望得到一段所指示的那种讲论的实例，好让我通过一个实际的例子向你展示我所推荐的事项如何执行。我确实将这样做，尽我所能藉天主的帮助。但在进行之前，按照我所承诺的，我有责任谈谈获得（我已提及的）喜乐的问题。因为关于你讲论应遵循的规则，在为一位明确表示愿意成为基督徒的人提供教理讲授的情况下，我已就我所承诺的内容，尽我所认为足够的程度作了说明。我自然没有义务同时在这卷书中亲自实践我所吩咐的正确做法。因此，如果我那样做了，那将是额外的好处。但在我完成应尽之责之前，怎能添加额外好处呢？此外，我听到你抱怨最多的一件事是，当你在教导某人基督徒的名号时，你的讲论似乎贫乏而无活力。我知道，这与其说是由于缺乏可说的内容——我深知你对此已有充分准备和装备——或言语本身的贫乏，不如说是由于心灵的厌倦。这或许源于我已经谈过的原因，即我们的理智更喜悦并更彻底地被我们静默地在心中所领悟的事物所吸引，而不愿将注意力从那里转移到远不能表达它的嘈杂言语上；或者由于即便讲论令人愉悦，我们更乐于聆听或阅读那些以更卓越的方式表达、且无需我们费心焦虑而呈现的内容，而非为了他人理解而拼凑出突然想到的言辞，其结果不确定：一方面，这些词语是否足以表达意思；另一方面，它们是否被接受而有益处；或者又由于这些内容对我们已完全熟悉，对我们自己的成长不再必要，因而频繁重复向未受过教导者宣讲的题材使我们生厌，而我们的心灵既已相当成熟，便毫无乐趣地在这类陈腐、近乎幼稚的题目上打转。当讲者面对一位无动于衷的听众时——他要么确实未被任何情感触动，要么未以任何身体动作表示理解或对所说内容感到满意——也会产生厌倦感。并非我们贪求人的称赞是合宜的态度，而是我们所施行的事属于天主；我们越爱那些听我们讲论的人，就越希望他们因听到为他们的得救而提出的内容而感到喜悦。因此，若我们未能做到这一点，我们便感到痛苦，力量减弱，中途灰心丧气，仿佛徒劳无功。有时，当我们被从某件我们渴望继续的工作（其进行令我们愉快，或显得格外必要）中拉走，并被迫——或因一位我们不愿冒犯之人的命令，或因一些无法摆脱之人的纠缠——去教导某人教理时，在这种情况下，我们带着已烦乱的心去接近一项需要极大平静的任务，既因我们不被允许按照自己所愿的顺序行事而悲伤，又因我们不可能胜任所有事情而痛苦；于是，讲论因出自沮丧的干涸土壤，自然变得不够流畅。有时，忧愁因某种冒犯占据了我们的心，恰恰那时有人对我们说：\Quote{来，与这人谈话；他愿意成为基督徒。} 因为这样对我们说话的人并不知道我们内心隐藏的困扰。于是，若他们不是我们应向其敞开情感的人，我们便毫无乐趣地承担他们所愿之事；那么，通过一颗既激动又焦躁的心所传出的讲论，必然是无力而乏味的。既然有如此多的原因遮蔽我们心灵的平静安宁，我们就必须按天主的旨意寻求补救之法，以缓解这些沉重之感，帮助我们以热切的精神喜乐，以善工的宁静欢愉。\BibleQuote{因为天主喜爱乐意捐献的人。}\BibleRef{格后9:7}\par

15. 如今，如果我们悲伤的原因在于，我们的听者未能理解我们的意思，以至于我们不得不以某种方式从自己概念的高度降下，被迫长久停留在远低于我们思想标准的冗长音节过程中，并且焦虑地思考，我们以最快的精神领悟所把握的东西，要如何通过肉体的口舌，以其复杂曲折的表达方式传达出来；而由此感受到的（言语与思想之间的）巨大差异，使我们不愿开口，却乐于沉默；那么，让我们思量那为我们树立了榜样、叫我们追随他足迹的那一位所展示的：\BibleQuote{给我们留下了榜样，叫我们追随他的足迹。}（见：\BibleRef{伯前 2:21}）因为无论我们清晰的言语与智力的活力有多么不同，可朽的肉体与天主的相等之间的差别还要大得多。然而，虽然\BibleQuote{他原具有天主的形体，却空虚了自己，取了奴仆的形体}，——直至\BibleQuote{十字架上的死亡}这些话语。（见：\BibleRef{斐 2:6-7}；\BibleRef{斐 2:8}）这对我们来说，除非他使自己成了\BibleQuote{对软弱的人，就成了软弱的，为赢得那软弱的人}（见：\BibleRef{格前 9:22}），否则还有什么解释呢？请听他的追随者在另一处表达类似意思的话：\BibleQuote{因为我们或是疯狂，那是为了天主；或是清醒，那是为了你们。原来基督的爱催迫着我们，因为我们曾如此断定：他替众人死了。}（见：\BibleRef{格后 5:13-14}）确实，如果一个人觉得俯就听众的耳朵是件烦人的事，又怎能准备好为他们灵魂的缘故而耗尽自己呢？因此，他就在我们中间成了一个小孩子，并且像保姆抚爱自己的孩子一样。因为，除非爱催促我们，否则用简短的、支离破碎的词语咿呀学语，难道是一种乐趣吗？然而人们却渴望有需要他们那样服侍的婴孩；母亲将一小块嚼碎的食物放进她小儿子嘴里，比她自己吃掉并吞下大块食物更为甜美。同样，因此，不要让母鸡的形象从你心中消失：她垂下羽毛遮盖自己柔弱的幼雏，并以断续的声音呼唤着她唧唧叫的孩子们到身边来，而那些因骄傲离开她呵护的翅膀的，就成了飞鸟的猎物。因为如果智慧在其最纯净的深处带来喜悦，那么我们也应当乐于智慧地理解这种方式：爱德越温顺地降卑到最卑微的对象，就越有力量沿着良心的路途回到最隐密的事物中，这良心见证她除了从所降卑之人身上寻求他们永远的救恩外，绝无他求。\par

% structure: p0031 source=h2 target=section reason=relative-heading-rank; base=h2

\section{第十一章 论第二厌倦根源的补救}

16. 然而，如果我们更愿阅读或聆听那些已为我们预备好并以优美风格表达的内容，结果便觉得自行组织言词既费时又结果不定，那么，只要我们的心思未偏离事实本身的真理，听者若因我们言辞中的任何不当而反感，他便很容易从这一情况中认识到：只要主题本身被正确理解，文字表达中即便有些缺陷或不准确，也不值得过分在意——因为这些表达本意仅是为确保对主题的正确领会。但若人性的软弱偏离了事实本身的真理——尽管在教导无知者的教理时，我们需遵循最常规的路径，这种情况不太容易发生——然而，为了避免听者在这方面也可能找到发怒的缘由，我们不应认为此事临到我们，除非是天主愿意借着它来考验我们，看我们是否愿意以平静的心接受纠正，以免在更大的错误中鲁莽地为我们自己的错误辩护。但若无人告诉我们这事，而它完全未引起我们自己或听者的注意，那么便不必为此忧伤，只要同样的事不再发生第二次。因为大多数情况下，当我们回想自己所说的话时，我们自己会发现一些可指责之处，却不知这些话在讲出时是如何被接受的；因此，当爱德在我们心中炽热时，若某件事实属错误却被无疑地接受，我们会更加烦恼。既然如此，每当机会出现——正如我们曾默默自责那样——我们也应当以体贴的方式纠正那些因我们的言辞（而非天主的圣言）而陷入某种错误的人。另一方面，若有人因愚昧的恶意而盲目欢喜我们犯了错误，他们是搬弄是非者、毁谤者、\BibleQuote{可恨天主的人}，这样的人应当成为我们以怜悯之心操练忍耐的材料，因为\BibleQuote{天主的忍耐引导他们悔改}。还有什么比效法魔鬼的邪恶形象，因他人的邪恶而欢喜，更可憎、更能在\BibleQuote{震怒与正义审判揭示之日}积聚震怒呢？有时，即使所有言辞正确且真实，也会因某句未被理解的话，或因与旧错误观念和习惯相悖而显得刺耳的新颖说法，使听者感到冒犯和不安。但若此事显露出来，且此人显示出可被纠正，便应立即用大量权威和理由纠正他。而若冒犯是隐晦的、不显明的，天主的良药便是有效的医治。再若此人退缩不肯被治愈，我们当以主耶稣的榜样安慰自己：当人们因祂的话而跌倒，视之为难懂的教训时，祂却对留下的人说：\BibleQuote{你们也愿意走吗？}我们心中应当坚守一个\Quote{坚定不动摇}的信念：被掳的耶路撒冷将在时期圆满时从这世界的巴比伦得到释放，属她的人没有一个会丧亡；因为凡丧亡的，本不属于她。\BibleQuote{天主的坚固根基屹立不移，且有这印记：主认识属于祂的人；凡称呼基督名号的人，应当远离不义。}我们若思虑这些，并呼求主进入我们心中，便会较少惧怕因听者不确定的感受而导致的言词结果；在为慈悲工作而忍耐烦扰时，烦扰本身也将令我们感到甘甜——只要我们在其中不寻求自己的荣耀。因为当行事者的动机为爱德所推动，并如回归本位般安息于爱德时，工作才真正是善的。此外，那令我们欣喜的阅读，或聆听高于我们自己的口才，其效果是使我们倾向于更珍视它胜过我们自己要发表的言论，从而令我们以勉强或冗长的言辞说话——但若能以更喜乐的心面对，并在辛劳后发现它更令人愉悦。再者，我们将更有信心向天主祈祷，愿祂按我们所愿对我们说话，因为我们欣然让祂按我们的能力藉我们说话。如此，万事互相效力，使爱天主的人获得益处。\par

% structure: p0033 source=h2 target=section reason=relative-heading-rank; base=h2

\section{第十二章——论第三种厌倦之源的补救}

17. 然而，我们常常觉得反覆讲述那些完全熟悉、且（更）适合儿童的内容是一件十分厌烦的事。如果我们自身有这种情况，我们就该尽力以兄弟、父亲和母亲的爱情来对待他们；一旦我们与他们心心相印，这些事物对我们和对他们都会显得新鲜。因为同情心态的力量如此巨大，以至于当他们听我们说话时深受感动，我们也在他们学习时深受感动，我们彼此居住在对方内；因此，在同一时刻，他们仿佛在我们内说出他们所听到的，我们也在他们内以某种方式学习我们所教导的。当我们向从未见过某些广阔美丽的地域（无论是在城市或田野）的人展示它们，而这些地域因我们习以为常而视若无睹、毫无愉悦时，我们岂不常发现，由于他们的新奇感，我们自己的享受也重新被点燃？而且，我们与他们的友谊越亲密，这种体验就越强烈；因为借着爱的纽带，我们存在于他们内，同样程度地，先前陈旧的事物重新变得新鲜。然而，如果我们自己在对事物的默观研究中取得了相当的进步，我们就不愿我们所爱的人仅仅因观赏人手所造的工程而感到满足和惊叹；而是愿意将他们提升到（默观）其作者的技巧或智慧本身，并由此（看到他们）升至对创造万有的天主的赞美和颂扬，在他那里是爱情最丰硕的终结。那么，当人们带着已经形成的、为获得天主本身的知识——而一切应学之事都应以此为目的——的心意来到我们面前时，我们岂不更应感到喜乐？我们岂不该因他们新颖的经验而更新自己，以至我们平常的讲道若有些冷淡，也会因他们非同寻常的听闻而重新燃起热情？此外，还有一点有助于我们达成喜乐，即我们默想并牢记此人正从错误的死亡过渡到信德的生命。如果我们怀着慈善的喜乐行走在最为熟悉的街道上，碰巧为一位因迷路而苦恼的人指路，那么，在救恩的道理上，我们岂不更应以高昂的精神和更大的喜乐，一遍又一遍地重新踏足那些就我们自己而言无需再开辟的路径？因为我们正带领一个可怜的灵魂，被这世界的歧途耗尽，行走在平安的道路上，顺服于那位为我们成就了这平安的天主！\par

% structure: p0035 source=h2 target=section reason=relative-heading-rank; base=h2

\section{第十三章——论第四种厌倦的补救方法}

18. 然而，说实话，当我们看到听者丝毫未被感动，却仍要继续讲说到预定限度，这实在是对我们的一个严峻要求；不论是因为他在宗教敬畏的约束下，不敢用声音或任何身体动作表示赞同，还是因为他被人的谦逊所抑制，或者他不理解我们所说的，或者他认为这些毫无价值。既然我们无法洞察他的心思，这一点对我们来说必定是不确定的，那么在我们的讲论中，我们有责任尝试一切可能有助于激发他、仿佛将他从隐藏之处引出来的方法。因为那种过度的、阻碍他表达判断的恐惧，应当被和善劝诫的力量驱散；通过向他提出我们兄弟般亲密关系的考虑，我们应当缓和他对我们的敬畏；通过询问他，我们应当确定他是否理解我们对他所说的话；我们应当向他传达一种自信感，使他能自由表达任何对他而言浮现的异议。同时，我们应当问他是否以前曾听过这类主题，以及是否因为这些内容对他而言如熟悉平常之事而未能感动他。我们应当根据他的回答调整我们的做法，要么以更简单的方式和更详细的解释来说话，要么驳斥某种敌对意见，要么，不对已知的主题进行更广泛的阐述，而是给予一个简短的总结，并选择一些在圣经中，尤其是在历史叙述中以奥秘方式处理的事，其展开和阐述可能使我们的讲话更吸引人。但如果此人天性极其迟钝，愚钝，与所有这类愉悦之源毫无共通之处，那么我们就必须以怜悯之心容忍他；在简要论及其他要点之后，我们应当以一种能激发他敬畏的方式，向他强调关于天主教会合一、关于试探、关于基督徒在未来的审判面前的生活等最不可或缺的真理；我们更应该为他向天主祈祷，而不是向他谈论许多关于天主的事。\par

19. 同样常见的是，起初满怀热忱听我们讲话的人，或因听讲的劳累，或因站立，变得疲惫，不再赞同所讲的内容，反而张口打哈欠，甚至不情愿地表现出想要离开的迹象。当我们察看到这一点，就有责任说一些融合了真诚愉悦、切合所讨论之主题的话，或说一些非常奇妙惊人、乃至可能是痛苦哀伤性质的话，以使他精神振奋。我们所说的话，若能更直接地触动他本人，使敏锐的自我关切之情令其注意力保持警觉，则效果更佳。然而，同时，所说的话也不应带有任何严厉而冒犯其虔敬之心；而应具有一种借所散发之友善来安抚他的性质。或者，我们应该让他坐下以缓解他的不适，尽管毫无疑问，如果从一开始，在适宜的情况下，他能坐着听讲，事情会处理得更好。事实上，在海外的一些教会中，出于对适宜性更周到的考虑，不仅主教们在向群众讲话时坐着，而且他们自己也为群众设置座位，以免任何体弱的人因站立而疲惫，因而其注意力从最有益的（讲道）宗旨上分散，甚至不得不离开。然而，如果只是人群中的某个人因体力不支而离开，且此人已因参与圣事而负有义务，则是一回事；但如果离开的人是（因为在此情况下，他通常不可避免地因担心自己会被内在的软弱压倒而跌倒，而被催促采取那种行动）将要受入门圣事的人，那又是另一回事；因为处于这种地位的人，一方面出于羞耻感而不便说出离开的理由，另一方面其体力又不允许他站立。这是我根据经验说的。因为在我以教理讲授的方式指导一位乡下来的人时，就曾发生过这种情况；从他那件事例中，我明白了这类事情必须谨慎防范。因为，如果我们不让那些是我们的弟兄、甚至那些尚未与我们处于这种关系的人（因为那时我们更应加倍关心，使他们成为我们的弟兄）坐在我们面前，谁能忍受我们的傲慢呢？更何况连一个女人坐着聆听我们的主亲自讲话，而天使们却警觉地侍立着！当然，如果讲道简短，或者场所不便于坐，他们应该站着听。但那只应在听众很多、且当时不准备正式接纳他们的情况下才如此。因为当听众只有一两个人或少数人，且他们明确为了成为基督徒而来时，站着对他们讲话是有风险的。不过，假设我们曾以那种方式开始，那么我们至少应在观察到听者疲惫迹象时，允许他坐下；不仅如此，我们还应尽力劝他坐下，并且应该说一些话，既能使他精神振作，又能驱散他心中可能偶然侵入并分散注意力的任何忧虑。因为既然我们无法确知他沉默不语、不愿听讲的原因，现在他既然坐下了，我们可以一定程度地针对世俗事务想法的侵入而发言，或以我已提及的愉快精神，或以严肃口吻；如此，若这些忧虑恰占据了其心，它们就会像被点名指控一般退去；而另一方面，若它们并非（失去兴趣的）特殊原因，而只是听讲疲劳，则当我们以我所说的方式，针对这些话题（仿佛它们就是那些特殊忧虑）发表一些出人意料、非同寻常的言论时，他的注意力就会从疲劳的压力下得到缓解——因为我们确实不知道（真正的原因）。然而，这样插入的话要简短，尤其考虑到它是额外插入的，以免这剂药反而加重我们想缓解的疲劳之症；同时，我们应该迅速继续未尽的论述，并承诺和呈现比预期更近的结论远景。\par

% structure: p0038 source=h2 target=section reason=relative-heading-rank; base=h2

\section{第十四章——针对第五和第六种倦怠根源的补救措施}

20. 再者，若你的心神因必须放弃某件你原以为更必要的其他工作而受挫，并因那种不得已而产生的忧伤使你带着不愉快的心情讲授要理，你应当思考这一点：除了我们知道在所有与人交往中，我们有责任以仁慈和至诚的爱德行事之外——我说的是这一例外——我们完全不确定何为更有益之事，何为更适时之事，无论是暂时搁置还是完全略过。因为我们不知道我们所代其功绩的人们在天主面前如何，所以关于某一时刻什么有益于他们的问题，我们非但不能理解，反而只能凭藉最微弱和最不确定的推断来猜测（甚至完全没有清晰的推断）。因此，我们当然应当按照自己的明智来处理我们的事务；然后，如能按决议的方式完成，就当欢喜，不是因我们的意愿，而是因天主的意愿使它们如此成就。但若发生任何不可避免的事，干扰了我们所提议的安排，我们应当顺从它，免得被折断；这样，天主所偏好而非我们所偏好的安排，也成为我们自己的。因为更合宜的是我们追随祂的旨意，而非祂追随我们的。此外，关于我们想按照自己的判断来保持的做事顺序，应当赞成那些优先考虑更高事物的方式。那么我们为何因主天主以其对我们的巨大优越性而凌驾于人之上而感到委屈，以至于在我们对自己顺序的爱中，竟如此轻视顺序？因为\Quote{没有人能更好地安排}他所要做的事，除非那人更愿意放下神能所禁止的事，而不是渴望做他以其人的思虑中所图谋的事。因为\Quote{人心有许多计谋，但主的计划却永远立定。}\par

21. 但若我们的心思因某种冒犯的原因而激动，以致不能进行平静而愉快的讲论，那么我们对基督为其而死的人的爱德，祂渴望用自己血的代价将他们从这世界错误的死亡中救赎出来——这爱德应当是如此巨大，以至于在我们悲伤时，那关于某个渴望成为基督徒的人到来的消息，就足以使我们欢喜，驱散心头的沉重，就像获利的喜悦通常能减轻损失之痛一样。因为我们并非因任何个别之人的冒犯而真正受压迫，除非是我们察觉或相信那人自己正在丧亡，或正成为某个软弱者跌倒的缘由。因此，一个人怀着被正式接纳的意愿来到我们这里，既然我们对他能前进抱有希望，就应该擦去因一个辜负我们的人所带来的悲伤。因为即使担心我们的皈依者会成为地狱之子的念头进入我们的思想——正如我们眼前有许多这样的人，那些让我们苦恼的冒犯正来自于他们——这不应阻碍我们，反而应激励和催促我们。同样，我们应当告诫我们所教导的人，要提防效法那些只在名义上是基督徒而非真正基督徒的人，并注意不要因他们人数众多而动摇，以致要么渴望跟从他们，要么因他们而不愿跟随基督，要么不愿意在天主的教会中与他们同在，要么希望成为他们那样的人。不知怎的，在这种告诫中，那由当下的悲伤感受提供燃料的讲论更加炽热；因此，我们在这种情绪下并非更迟钝，而是以更大的火热和激烈说出那些在更安逸时我们会以更冷淡和更无力方式说出的事情。这应当使我们欢喜，因为有一个机会使得我们内心的情感不会毫无成果地消逝。\par

然而，若因我们自己曾犯差误或罪过而陷入忧伤，我们不仅应记住\Quote{心灵破碎是向天主的祭献}，也当记住\Quote{水能扑灭火，施舍能赎罪}；祂又说：\Quote{我喜欢仁慈，胜过祭献}。因此，就如当我们面临火灾危险时，必定会跑向水去灭火，且若有人就近提供水，我们必会感激；同样，若从我们自己的堆中升起罪的火焰，且我们因此感到困扰，当有了一个极富仁慈工作的机会时，我们就应为此喜乐，仿佛有一道泉水被提供给我们，好使已爆发的烈火藉以熄灭。除非我们愚蠢到认为，我们应更乐于带着面包跑去填满饥饿之人的肚腹，而胜过带着天主的言语去教导那以此喂养者的心灵。这一点也值得考虑：即倘若做此事只对我们有益，而放弃却无损害，我们或许会轻视一种在危险时刻以不幸方式提供的补救，这补救关乎的，不是邻人的，而是我们自己的安全。但当我们从主的口中听到如此威胁性的判词：\Quote{你这可恶而又懒惰的仆人！你该把我的金钱交给兑换银钱的人}，请问，既然我们的罪使我们痛苦，那么因拒绝将主的金钱给那渴望和请求的人，而决意再次犯罪，这是何等疯狂！当这些以及类似的思考和反思成功驱散了倦怠情绪之黑暗后，心灵的方向便适于训导的职责，从而使那从丰富的爱德脉络中活泼而欢快地涌出的话语，以愉悦的方式被接受。因为这些在此所说的话，与其说是由我对你们说的，不如说是由那\Quote{藉着赐予我们的圣神，倾注在我们心中的爱}对我们众人说的。\par

% structure: p0042 source=h2 target=section reason=relative-heading-rank; base=h2

\section{第十五章　论应如何使我们的训导适应不同类别的听众}

23. 但如今，或许你也以债的方式向我索求那在做出承诺之前本无义务给予的东西，即我不应视为累赘地展开某种意向中的论述范例，并将其呈于你面前供你研习，就仿佛我正在亲自教导某个人一样。然而在这样做之前，我愿你记住：一个人心中想着未来的读者而口述时，心思投入是一种；而当他对着眼前的听者说话并留意其反应时，则是另一种。此外，尤须牢记：在后一种情况中，当私下一对一地劝诫时，且没有旁人在场评判我们时，心思投入是一种；而当公开传授教导，且周围站着意见各异的听众时，则是另一种。再者，在此教导实践中，当只有一人受教，其余人则像评判或见证已知之事那样聆听时，心思投入是一种；而当所有在场者都期待我们向他们宣讲时，则是另一种。同样地，在同一个场合中，当所有人仿佛私密会议般围坐以进行讨论时，心思投入是一种；而当人们静默地坐着，全神贯注于一位从较高位置向他们说话的讲者时，则是另一种。同样地，即使我们以那种风格讲论，听众的多少、是否有学识、是纯然一种还是两者混合、是城里人还是乡下人、还是二者都有、或者是由各阶层按比例组成的人群，这些都会造成显著差异。因为这些人必然以不同方式影响必须对他们说话和讲论的人，而所发表的讲话既会带有某种特征，仿佛表达出产生它的心灵情感，也会根据（讲者心境的）同样差异以不同方式影响听者，同时听者之间也因彼此在场而相互影响。但既然目前我们讨论的是教导无知者的事，我以自己的经历向你作证：我发现当我在教理讲授中面对不同的人时——如博学之士、愚钝之人、市民、外国人、富人、穷人、平民、显贵、有权势者、不同民族、不同年龄或性别的人、来自不同派别或不同一般错误的人——我内心的感动是多种多样的，而我的论述也总是随我心境的差异而开始、进行并结束。而且，虽然同一爱德当施于众人，但同一药物不可施于众人；同样，爱德本身与某些人一同劳作，与另一些人一同软弱；费力建立一些人，战战兢兢怕得罪一些人；向一些人屈身，向另一些人挺立；对一些人温和，对另一些人严厉；对无人是敌人，对众人是母亲。当一个并未以同样爱德精神经历我所提及的这类体验的人，看见我们因蒙受某种赐予的恩赐——此恩赐带有使人喜悦的能力——而在众人口中获得某种赞美性的声誉时，他便会因此视我们为有福。但愿那鉴察\BibleQuote{被囚之人的叹息}的天主，垂视我们的谦卑和劳苦，并宽恕我们所有的罪过。因此，如果我们身上有任何事物至今令你喜悦，以致你渴望听我们谈论关于你当遵循的论述形式的意见，你应当通过观察我们、听我们实际讨论这些主题时的情形，来获得对这一问题更深入的理解，胜过通过阅读我们口述时的文字。\par

% structure: p0044 source=h2 target=section reason=relative-heading-rank; base=h2

\section{第十六章 教理讲授范例；首先，一个心怀合适宗旨的慕道者的案例}

24. 尽管如此，我们不妨假设有人来到我们这里，希望成为基督徒；他并非出身乡野，而是来自城市，就像你在迦太基必定会遇到的许多人那样。我们也假设，当被问及他渴望成为基督徒的原因，是基于现世的好处，还是寄望于来世的安息时，他回答说，他的动机是来世的安息。那么，这样一个人或许可以由我们以下列方式教导：\Quote{感谢天主，我的弟兄；我由衷地为你高兴，并因你在这既大又危险的世间风浪中，想到了某种真实而稳固的保障而欣喜。因为即便在今世，人们也以沉重的劳苦寻求安息和安稳，但因他们的邪恶私欲而找不到。他们以为在不安静且不持久的事物中能找到安息。这些事物，由于从他们那里被夺去并在时间中消逝，以恐惧和忧伤搅扰他们，不容他们享受宁静。因为如果一个人想在财富中寻找安息，他会变得傲慢而不是平静。我们岂不见多少人突然失去财富？多少人因财富而毁灭，或因贪图拥有，或被更贪婪的人压倒并夺去？即使财富终生陪伴他，永不离开它的爱慕者，他本人也必在死亡时离开它们。人的生命，即使活到老年，又有多大呢？或者，当人们为自己渴望老年时，他们实际上渴望的不过是长久的虚弱。同样，世上的荣华——不过是空虚的骄傲和虚荣，以及堕落的危险。因为圣经这样说：\BibleQuote{凡有血肉的，尽都如草，他的美容，如同草上的花。草必枯萎，花必凋谢；但上主的话永远存留。}\BibleRef{依 40:6-8} 因此，如果有人向往真正的安息和真正的幸福，他应将希望从必死和短暂的事物上挪开，并固定在主的话上；这样，因依附那永久长存者，他自己也与他一同永久长存。}\par

25. 还有另外一些人，他们既不想发财，也不追求虚浮的荣华，却只愿在美味佳肴、奸淫以及大城市中免费提供的剧场和表演中寻求快乐和安息。但他们的遭遇也是如此；他们要么在小事上浪费微薄的资产，随后因贫乏而行窃、入室抢劫，甚至拦路抢劫，于是突然陷入众多且巨大的恐惧中；刚才还在欢宴之家歌唱的人，现在却梦想着监狱的悲伤。此外，他们热衷公共表演，变得如同魔鬼，因为他们以呼喊激励人们互相伤害，煽动那些未曾伤害过他们的人彼此激烈争斗，同时却要取悦疯狂的人民。如果他们发现任何人和平相处，便立刻仇恨并迫害他们，并大声叫嚣，要求用棍棒殴打他们，仿佛他们勾结欺骗了自己；这种不义甚至迫使法官——本该是惩罚不义的——去施行。另一方面，如果他们看到那些人彼此进行可怕的敌对，无论他们是被称为\Emph{sintœ}，还是戏剧演员或表演者，或是赛车手，或是猎人——那些他们使之争斗并较量（不仅是人与人，甚至是人与兽）的不幸之人——那么，他们看到这些可怜的人彼此愤怒地冲突得越猛烈，就越喜欢他们，越从中获得乐趣；他们帮助这些激动的人，而通过帮助他们，越发激发他们；观众们彼此争斗，支持不同的斗士，甚至比那些他们疯狂激怒以致疯狂的人还要疯狂，同时他们自己也在疯狂中渴望成为同样的观众。那么，那以纷争和争斗为食的心思，怎能保持和平的健康呢？因为正如所接受的食粮，所产生的健康也是如此。总之，尽管疯狂的快乐并非快乐，但我们且就事论事：无论这些事物的性质如何，也无论它们所带来的快乐程度如何——财富的夸耀、荣华的膨胀、酒馆的挥霍享乐、剧场的争斗、奸淫的污秽、浴池的淫欲——所有这些都是小小热病就能夺走的，甚至对于那些还活着的人，也夺走了他们生活中所有虚假的快乐。剩下的只有空虚和受伤的良心，注定要视那他不愿认作父亲的天主为审判者，也注定要在他不屑寻求和爱慕的温柔父亲身上，找到一个严厉的主。但你，既然寻求那在今生之后应许给基督徒的真正安息，即使在此生最痛苦的患难中，也必尝到同样甜蜜而愉快的安息，只要你继续爱那应许者的诫命。因为你将很快感到正义的果实比不义的果实更甜美，并且一个人在患难中拥有良心的平安，比在快乐中拥有邪恶的良心更真实、更愉快。因为你来到天主的教会，并非为了从中寻求任何现世的利益。\par

% structure: p0047 source=h2 target=section reason=relative-heading-rank; base=h2

\section{第十七章——继续教理讲授范例，特别针对斥责慕道者方面的虚假目标}

26. 因为有些人渴望成为基督徒的理由，要么是为了从他们期待获得现世利益的人那里得到宠爱，要么是不愿得罪他们所畏惧的人。但这些都是被弃绝的人；虽然教会暂时包容他们，就像打谷场包容糠秕直到簸扬的时候，但若他们不改正，不以真诚的心成为基督徒，着眼于将要来临的永恒安息，最终他们必从教会中被分离出去。愿这样的人不要自欺，因为他们有可能与天主的谷粒同在打谷场上。但他们不会与谷粒一同进入谷仓，而是注定被投入火中，那是他们应得的。还有另一些人，固然有更好的希望，但同样处于不小的危险中。我指的是那些如今敬畏天主，不嘲笑基督徒的名号，也不是以虚伪的心进入天主的教会，但仍在此生寻求幸福，期望在尘世事物中比那些拒绝敬拜天主的人享有更多福乐的人。这种虚假期望的后果是，当他们看到某些邪恶不虔之人权势稳固、在今世亨通，而自己要么拥有的更少，要么完全失去时，便因思虑自己事奉天主没有理由而困扰，于是轻易从信德中堕落。\par

27. 至于那着眼于今世之后的永恒福乐和永久安息（这是应许给圣人们作为命运的份），并渴望成为基督徒，以免与魔鬼一同进入永火，而与基督一同进入永恒国度的人，这样的人是真正的基督徒；（并且他）在一切诱惑中保持警醒，既不因顺境而败坏，也不因逆境而完全沮丧，而是当世上的好事丰富时保持谦逊节制，当患难临到时勇敢坚忍。具有这种品格的人也将不断进步，直到达到那种心态：他爱天主胜过畏惧地狱；即使天主对他说：\Quote{永远享受肉体的快乐，尽情犯罪吧，你不会死，也不会被投入地狱，但你只是不能与我同在}，他也会极其惊骇，完全戒绝犯罪，现在不再（仅仅）是为了不落入他所害怕的那种境况，而是为了不冒犯他如此深爱的那一位；唯有在祂内，也有那\BibleQuote{眼所未见，耳所未闻，人心所未曾想到的——就是天主为爱祂的人所预备的安息。}\par

28. 关于这安息，圣经是意味深长的，且不缄默地讲述，它告诉我们世界之初，天主创造天地和其中万物的那时，祂工作了六天，并在第七天安息了。因为全能者本可以在瞬间创造万物。祂并不劳苦，以致需要安息，因为的确\BibleQuote{祂一发言，万有造成；祂一命令，万物生成。}但祂要表明，在世界六个时代之后，在第七个时代，如在第七天，祂要在祂的圣人们内安息；因为这些圣人们也要在祂内安息，在他们为祂所做的一切善工之后——这些善工实际上是祂自己在他们内完成的，祂召叫他们，教导他们，除去过去的过犯，并使先前不虔的人成义。因为当祂的恩赐使他们行善时，正确地说是祂自己（在他们内）工作；同样，当他们在祂内安息时，正确地说是祂自己安息。至于祂本身，祂不寻求休止，因为祂不感到劳苦。此外，祂藉着祂的圣言创造万有；而祂的圣言就是基督自己，天使和天上所有最纯洁的灵体都在神圣的静默中安息于祂内。然而人，因着罪恶而堕落，失去了他在天主性中所拥有的安息，现在在祂的人性中重新获得它；为此目的，祂成为人，由女人所生，在祂自己知道应当如此完成的适宜时刻。从肉体祂确实不会受到任何污染，因为祂反而注定要净化肉体。古代圣人们，在圣神的启示中，预知了祂的来临，并预言了。因此他们因相信祂将要来临而得救，正如我们因相信祂已经来临而得救。因此我们应当爱天主，祂如此爱了我们，以至派遣了祂的独生圣子，使祂披上我们必死性的卑微，并死于罪人之手，也为罪人而死。因为在古代，在最初的世代，这奥秘的深度不断被预表和预言所宣告。\par

% structure: p0051 source=h2 target=section reason=relative-heading-rank; base=h2

\section{第十八章——关于人类及受造物的创造所当信的内容}

29. 既然全能的、美善、公义而仁慈的天主，创造了万物——无论大小、高低，无论可见的，如天、地、海，以及天上特别的太阳、月亮和其他发光体，以及地上和海中的树木、灌木、各类动物，以及一切天上地下的物体，或不可见的，如那些使躯体具有生命和活力的精神体——也按自己的肖像造了人，以便正如祂自己因其全能而掌管整个受造界，人也能凭其智慧（借以认识并崇拜其造物主）掌管地上所有活物；祂又造了女人作为他的\OriginalTerm{助手}{helpmeet}，不是为了肉身的\OriginalTerm{私欲偏情}{concupiscence}（因为在死亡形式的罪的惩罚侵入之前，他们那时并无可朽坏的身体），而是为了这个目的：男人可因女人而荣耀，因为他在走向天主方面超越她，并以自身作为圣洁和虔敬的榜样供她效法，正如他自己因追随智慧而成为天主的荣耀。\par

30. 因此，祂将他们安置在一个永福的地方，圣经称之为\OriginalTerm{乐园}{Paradise}，并给了他们一条诫命，若他们不违反，就永远留在那不朽的福乐中；若他们违反，则承受死亡的惩罚。天主早已预知他们会违反。然而，因为祂是一切美善的作者和创造者，祂宁愿造他们，如同祂也造了野兽一样，以便用地上固有的美善充满大地。而且，人，即使是有罪的人，也优于野兽。祂仍愿意给出那条他们不会遵守的诫命，为使他们当祂开始对他们施行惩罚时无可推诿。因为无论人做了什么，他都会发现天主在一切作为中堪受赞美：若他行为正直，他必因祂赏报的公义而赞美祂；若他犯了罪，他必因祂惩罚的公义而赞美祂；若他认罪并转向正直生活，他必因祂赦罪恩典的仁慈而赞美祂。那么，天主为何不应造人，尽管祂预知他会犯罪？因为祂可以在人站立时赏报他，跌倒时扶正他，站起时帮助他，而祂自己总是无处不在地在美善、公义和仁慈中得荣耀。尤其是，祂为何不应这样做，因为祂也预知这一点：从那有死的人类中，将有圣人兴起，他们不求自己的益处，而是归光荣于他们的造物主；他们借崇拜祂而获得摆脱一切腐败的拯救，堪当永远活着，并与圣天使一同在福乐中生活。因为祂将\OriginalTerm{自由意志}{free will}赐予人，为使人不是出于奴役性的必然，而是出于\OriginalTerm{真诚的倾向}{ingenuous inclination}崇拜天主；祂也将其赐予天使；因此，那位天使，与其他作为其部下的精神体一起，因骄傲背离对天主的服从而成为魔鬼，并未损害天主，而是损害了自己。因为天主知道如何处理那些离开祂的灵魂，并从他们公义的悲惨中，为祂受造物的较低部分提供祂奇妙\OriginalTerm{安排}{dispensation}中最恰当合适的法则。因此，魔鬼无论在自己堕落时，还是在诱惑人陷入死亡时，都丝毫未损及天主；人自己也一点没有损害造物主的真理、权能和福乐，因为当他的伴侣被魔鬼诱惑时，他出于自己故意的倾向赞同了她，去行天主所禁止的事。因为天主的公义法律定了所有人的罪，天主自己在报应的公义中得荣耀，而他们因惩罚的贬抑而蒙羞；为使人背离造物主时被魔鬼战胜并成为其奴仆，并使魔鬼作为必须被征服的敌人摆在人面前，当人重新归向造物主时；这样，凡至终同意魔鬼的，便与他同受永罚；而那些在天主前谦卑自己，并靠祂的恩宠战胜魔鬼的，堪受永赏。\par

% structure: p0054 source=h2 target=section reason=relative-heading-rank; base=h2

\section{第19章　论教会内善与恶的共存及其最终分离}

31. 也不要因许多人随从魔鬼、少数人跟随天主这一事实而动摇；因为麦子与糠秕相比，在数量上也大为欠缺。但正如农夫知道如何处理那大堆的糠秕一样，天主的眼里罪人的众多也毫无意义，祂知道如何处理他们，以致不让祂王国的管理在任何部分失序或蒙羞。也不应以为魔鬼因引诱了比少数能胜过他的人更多的人而就是胜利者。这样，有两个团体——一个是不敬神者的，另一个是圣者的——它们从人类之初一直延续到世界终了，目前它们在身体上混杂在一起，但在意志上却是分离的，并且，在审判之日，它们在身体上也将被分离。因为所有以虚妄的骄矜和傲慢的浮华爱慕骄傲和世俗权力的人，以及所有倾心于这些事物、在臣服他人中寻求自己光荣的邪灵，都紧紧联结在一个联盟中；即使他们因这些事常常互相争斗，他们仍被同样的贪欲之重压抛入同一个深渊，并且因相似的习性和功过而彼此联合。再者，所有谦卑地寻求天主的荣耀而非自己的荣耀，并在虔诚中跟随祂的人和邪灵，都属于同一个团体。尽管如此，天主对不敬神的人极其仁慈和忍耐，赐予他们忏悔和改过的空间。\par

32. 因为关于洪水毁灭所有人，只留下一个义人及其全家（祂愿意他们在方舟中得救）这件事，祂确实知道他们不会改正自己；然而，在建造方舟的一百年间，将要临到他们身上的天主义怒确实曾被宣讲给他们；如果他们愿意归向天主，祂就会宽恕他们，正如后来祂因尼尼微城忏悔而宽恕了它那样，此前祂曾通过一位先知向他们宣布即将临到的毁灭。同样，天主行事，即使对祂知道将坚持作恶的人，也赐予忏悔的空间，好藉祂自己的榜样操练并教导我们的忍耐；这样，我们也就能知道，当我们不知道他们将来会成为怎样的人时，用恒久忍耐宽容恶人是多么合宜，因为祂的认知无任何未成就之事可逃避，祂仍宽容他们，让他们活着。然而，在洪水的圣事性标记中，义人藉木头获救，这也预先宣告了将要形成的教会，即基督——祂的君王和天主——藉祂十字架的奥迹，将她高举于这世界的淹没之外。此外，天主并非不知道，即使那些在方舟里得救的人，也会生出恶人，再次以不义充满地面。然而，祂仍然赐给了他们将来审判的预像，并藉木头的奥迹预先宣告了圣者的解救。因为在这些事之后，邪恶并未停止通过骄傲、情欲和非法的不敬虔再次萌发，那时人们离弃造物主，不但堕落到所造的受造物（的标准）去敬拜受造物而非天主，甚至将他们的灵魂屈从于人的手工艺品和工匠的发明，魔鬼和那些邪恶的邪灵在这些虚假的伎俩中因被尊崇和敬拜而欣喜，同时用人的错误喂养自己的错误，从而在他们身上赢得了更可耻的胜利。\par

33. 然而，事实上，在那个时代也不乏义人，他们虔诚寻求天主，并战胜了魔鬼的骄傲，他们是那神圣团体的公民，藉着基督的谦卑而得痊愈；这谦卑当时只是预定要来临，但已藉圣神启示给他们。从这些人中，拣选了亚巴郎，一位虔诚而忠信的天主仆人，为向他显示天主圣子的圣事，好使将来万邦中所有信者因效法他的信德而被称为他的子女。从亚巴郎生出一个民族，当其他民族都事奉偶像和邪神时，这个民族应敬拜那创造天地的唯一真天主。显然，在这个民族中，未来的教会得到更清楚的预表。因为其中有属肉体的群众，为现世的利益而敬拜天主。但其中也有少数人思想未来的安息，切望天上的祖国；通过预言，他们得知我们的君王和主耶稣基督的天主降生成人的谦卑，为使他们藉此信德从一切骄傲和傲慢中获得痊愈。至于这些在时间上先于主诞生而存在的圣人们，不仅他们的言语，就连他们的生活、婚姻、子女和行为，都成为这个时代的预言，此时教会正藉着信德、通过基督的苦难从万邦中聚集起来。藉着这些圣祖和先知们，这个属肉体的以色列民族（后来也称为犹太人）同时获得了他们以肉体方式急切向主祈求的可见恩惠，以及旨在以肉体惩罚的形式暂时惊吓他们的惩戒，这适合他们的顽固。然而，在这些事中，也象征着精神的奥秘，这些奥秘是指向基督和教会的；那些圣人也是这教会的成员，尽管他们在主基督按肉体出生之前已存在于今生。因为这位基督，天主的独生子，圣父的圣言，与圣父同等、同永恒，万物藉祂而造成的，祂自己为了我们而成为人，为能作整个教会——如同祂整个身体——的元首。正如当整个人在诞生过程中，虽然手可能在出生时先伸出，但那只手与整个身体一起在元首之下连接并紧密结合，同样，在这些圣祖中，有些人在出生时手先伸出，作为这事的预兆：因此，所有在我们主耶稣基督诞生之前生活在世上的圣人们，虽然他们先出生，但都在元首之下与那普遍性的身体——祂是这身体的元首——联合在一起。\par

% structure: p0058 source=h2 target=section reason=relative-heading-rank; base=h2

\section{第二十章　论以色列在埃及为奴、他们的得救以及过红海}

34. 那时，那民族被带到埃及，受最残暴的国王奴役；因极沉重的苦役而被教导，他们就在天主面前寻求拯救者；于是从本族中给他们派遣了梅瑟，天主的圣仆，他藉天主的权能，用伟大的奇迹恐吓了当时埃及的亵渎民族，并带领天主的子民离开那地，经过红海；海水分开，为他们开路渡海，而当埃及人追赶时，波浪流回河床，淹没他们，使他们灭亡。这样，正如大地藉洪水从罪人的邪恶中被水洁净，那些时代的人被洪水毁灭，而义人藉木头得救；同样，天主子民离开埃及时，在水中找到一条路，而他们的敌人被水吞没。在那里也不缺少木头的圣事，因为梅瑟用他的杖击打，才成就了那奇迹。这两者都是神圣圣洗的标记，信者藉此进入新生，而他们的罪如同仇敌被消灭。但在那民族中，基督的苦难更清楚地被预表，当时他们奉命宰杀并吃羔羊，用它的血涂在门框上，每年举行这仪式，称之为主的逾越节。的确，预言极其清楚地讲论主耶稣基督，当它说\Quote{祂被牵去，像一只待宰的羔羊。}而今天，你要像在门框上一样，在你的额头上印上祂苦难和十字架的标记，所有基督徒都印有同样的标记。\par

35. 此后，这民族在旷野中行走了四十年。他们也领受了由天主手指所写的法律；在福音中最明确地宣示，这名称所指的乃是圣神。因为天主不以身体的形状来界定，我们也绝不可认为祂有肢体和手指，如同我们在自己身上所看到的那样。而是藉着圣神，天主的恩赐分给祂的圣人，使他们虽能力各异，却不致脱离爱德的和谐；又因在手指上特别显出一种分别，却与统一并不分离，这或许是这一说法的解释。但无论是否如此，或可能有其他理由将圣神称为天主的手指，我们听闻此语时，总不可思想人身的形状。这民族领受了由天主手指所写的法律，确实是写在石版上，以象征他们心硬而不能履行法律。因为他们急切向主求取身体使用的恩赐，被肉性的畏惧而非属神的爱德所束缚。但惟独爱德能成全法律。因此，他们被许多可见的圣事所重压，为使他们感受到奴役之轭的压力，即饮食的条例、牲畜的祭献以及其他无数的礼仪；这些同时又是与主耶稣基督和教会有关的属神之事的标记；此外，当时有少数圣人也领悟其义，以结出救恩的果实，并按时代的适宜而持守它们，而多数属肉性的人只是持守，却不领悟。\par

36. 这样，藉着许多预象未来之事的不同标记——若详尽列举将使人厌倦，且如今我们在教会中看到它们实现——那民族被领到了预许之地，要在那里按他们的渴望，以暂时和肉性的方式为王：这地上的国度，仍维持着属神国度的形象。那里建立了耶路撒冷，那座最著名的天主之城；它在奴役中，却预表了自由之城，即被称为天上耶路撒冷的城；这后一名称是希伯来语，解释为\Quote{和平的异象}。其所有公民都是成圣的人——无论是已经存在的、现今存在的、还是将要存在的；以及所有成圣的神体，就是那些在天上崇高之处以虔诚的敬爱服从天主、不效法魔鬼及其天使不敬虔骄傲的众神体。这城的君王是主耶稣基督，天主的圣言，最高天使由祂治理，同时祂又是取了人性的圣言，为使人也由祂治理，凡与祂共融的人，将在永远的和平中一同为王。为预表这君王，在以色列民族的地上国度中，达味王卓然显出，按肉身而论，那最真实的君王出自他的后裔，即我们的主耶稣基督，\Quote{祂是万有之上，永远受赞美的天主}。在预许之地，发生了许多事，都是未来基督和教会的预像；你将有机会在圣经中逐渐认识它们。\par

% structure: p0062 source=h2 target=section reason=relative-heading-rank; base=h2

\section{论巴比伦囚掳及其所预表之事}

37. 然而，经过几代之后，出现了另一个预象，与当前议题密切相关。因为那座城被掳掠，大批人民被带往巴比伦。如今，正如耶路撒冷象征圣人的城邦与团体，巴比伦则象征恶人的城邦与团体，因为按解释它意指\OriginalTerm{混乱}{confusion}。关于这两座城，它们自人类之初便交织运行，历经时代变迁，并将一同延续直到世界终结，那时它们将在最后审判中被分开——我们不久前已有所论述。当时耶路撒冷城被掳，人民被奴役地带往巴比伦，这一切是由主通过当时的先知耶肋米亚的声音所命定的。出现了巴比伦诸王，人民在他们手下为奴；由于这民族被掳的机缘，他们因某些奇迹而深受感动，以致认识了创造整个宇宙的唯一真天主，并敬拜祂，且命令他人也当敬拜祂。此外，人民奉命为俘虏他们的人祈祷，并在他们的和平中盼望和平，以至于生养儿女、建造房屋、种植葡萄园和果园。而在七十年结束时，他们获许诺得释放。这一切更以预象指明，基督的教会——所有属于天上耶路撒冷的圣人——将在这世界的君王下服务。因为宗徒的训导也这样说：\BibleQuote{人人都当服从更高的权力}，并且\BibleQuote{当向众人提供一切，向谁纳税就纳税，向谁纳粮就纳粮}，以及类似的其他事情，我们在不损害对天主的敬拜的前提下，向人类社会中的统治者提供这些；因为主自己，为了给我们树立这健全教义的榜样，也不认为因祂所取的人性而纳税是可耻的。同样，基督徒仆人和良好信徒也被命令以平和与忠诚侍奉他们现世的主人；如果最终发现他们是恶人，他们将审判他们；或者如果他们自己也被归向真天主，他们将与他们平等地统治。但所有人都被命令要服从属人和属世的权力，直到预定期限结束（这由七十年所预表），那时教会将从这世界的混乱中得释放，正如耶路撒冷将从巴比伦的掳掠中得自由。然而，借着那掳掠的机缘，地上的诸王自己也已被引导放弃那些他们曾用以迫害基督徒的偶像，并已认识并如今敬拜唯一的真天主和基督、主；宗徒保禄也命令为他们祈祷，即使他们迫害教会。因为他这样说：\BibleQuote{因此我首先恳求，要为所有君王、为所有人、为所有掌权者，献上恳求、祈祷、转求和感恩，使我们能虔敬和爱德地度安静和平的生活。}因此，教会从这些人得到了和平，虽是暂时的和平——一种暂时的安宁，用于以属灵的方式建造房屋和种植葡萄园与果园。因为请看你们自己的情形，就在此时，我们正借着这篇训导，在建造和种植你们。同样，借着基督徒君王所赐的和平，在整个世界范围内也进行着类似的过程，正如同一宗徒所说：\BibleQuote{你们是天主的田地，是天主的建筑物。}\par

38. 事实上，在耶肋米亚所预言、为预象征时期终结的七十年过去之后，为了完善那个预象，犹太人并未获得稳固的和平与自由。因此，他们随后被罗马人征服，成为进贡者。从他们获得应许之地并开始有君王的时候起，为了排除那将要来的、作为他们解放者的基督的应许在他们任何一位君王身上已完全实现的猜测，基督在诸多预言中被更清楚地预告：不仅由达味自己在圣咏集中，也由其他伟大圣洁的先知们，直到他们被掳到巴比伦之时；而在那同一掳掠中，也有先知奉命预告主耶稣基督作为所有人的解放者的来临。当圣殿在七十年后重建时，犹太人受到外邦君王的沉重压迫和苦难，这使他们明白解放者尚未到来——他们未能领悟祂是为他们带来属灵解放的，而他们却因属肉体的解放而徒然渴望祂。\par

% structure: p0065 source=h2 target=section reason=relative-heading-rank; base=h2

\section{第二十二章 论世界的六个时代}

39. 据此，世界的五个时代现已完成，（第六个时代已经进入。）这些时代中，第一个是从人类起源，即从第一个受造的人\Person{亚当}，直到洪水时期建造方舟的\Person{诺亚}。第二个则是从那个时期直到\Person{亚伯拉罕}，他被称为所有效法他信德的民族的父亲，但同时也按肉身的自然后裔，成为预定选民犹太人的父亲；在外邦人进入基督信仰之前，这选民是普世万民中唯一崇拜独一真天主的民族；\Person{基督救主}也按肉身预定从这民族而来。因为这两个时代的转折点在古经中占有显著地位。另一方面，其他三个时代也在福音中宣告，那里也提及了\Person{主耶稣基督}按肉身的谱系。因为第三个时代从\Person{亚伯拉罕}直到\Person{达味}君王；第四个时代从\Person{达味}直到天主子民被掳往巴比伦的那次充军；第五个时代从那次流徙直到我主\Person{耶稣基督}的来临。随着祂的来临，第六个时代已开始其进程；如此，那在先前只为少数宗祖和先知所知的神性恩宠，如今要向万民彰显；为的是叫人不可出于私欲崇拜天主，贪图祂服务的可见报酬和现世幸福，而只渴求那永恒的生命，在其中享受天主本身：为的是在这第六时代中，人的心灵能按天主的肖像得到更新，如同第六日人按天主的肖像被造。因为那时，法律也得以圆满，即凡它所命令的，都非出于对暂时之物的强烈欲望，而是出于对颁布命令者的爱。此外，有谁不应热切地以爱回报至上公义又至仁慈的天主呢？祂先爱了那些最不义、最傲慢的人，且爱得如此深，竟为了他们派遣了祂的独生子，万物由祂而造，祂成为人，并非出于祂自身的改变，而是因取了人性，为的是不仅能够与他们一同生活，而且还要同时为他们并借他们的手而死？\par

40. 如此，在彰显我们永恒产业的新约时，人将要因天主的恩宠而更新，并度一种新生活，即属神的生活；并且为了表明第一部约是旧约的安排，其中属血肉的百姓（少数有领悟力的宗祖、先知和一些隐藏的圣人除外）度着属血肉的生活，行事为人如同旧人，期望从上主天主那里得到属血肉的报酬，并这样接受了属神恩惠的预象；——怀着这个意向，我主基督成为人时，轻视了世上一切美物，为向我们显示这些美物应当被轻视；祂忍受了世上一切灾祸，并教导这些灾祸必须忍受；这样，我们的幸福不在前者中寻求，我们的不幸也不在后者中恐惧。因为祂由一位母亲所生，虽然她未与男子接触便怀孕，并始终如此，童贞受孕、童贞生产、童贞去世，却许配给一位工匠；祂消灭了属血肉的高贵的一切虚骄。再者，祂生于\Place{白冷城}，在犹大各城中如此微不足道，以致于今天仍被称为村庄；祂不愿有人因地上任何城市的显赫地位而夸耀。祂，万有属于祂，万物由祂所造，却成为贫穷的，为使人不得因信祂而敢于在属世的财富中自夸。祂拒绝被人立为王，因为祂向那些被骄傲从祂分离的不幸者展示了谦逊的道路；然而整个受造界都证明祂永恒王国的真实。祂使众人得饱，自己却饥饿；凡饮的都出于祂所造，祂是饥饿者的食粮和口渴者的泉源，自己却口渴；祂在地上旅行，疲惫不堪，却使自己成为我们进入天堂的道路；在辱骂祂的人面前，如同又聋又哑的人，而祂使哑者说话，聋者听见；祂被捆绑，却从软弱的捆绑中释放了我们；祂被鞭打，却从人的身体中驱除了一切痛苦的鞭打；祂被钉十字架，却终结了我们的剧痛；祂使死人复活，自己却死了。但祂也复活了，不再死亡，为使人不得从祂学习如此轻视死亡，仿佛自己永不再活似的。\par

% structure: p0068 source=h2 target=section reason=relative-heading-rank; base=h2

\section{二十三、论圣神在基督复活后五十天的降临}

41. 此后，耶稣坚固了门徒，并与他们同住了四十天，就在他们目睹之下升天而去。从祂复活算起，满了五十天，祂派遣了圣神给他们（正如祂所应许的），藉着圣神的德能，他们将爱德倾注在心中，从而能够遵守法律，不仅不觉得沉重，反而满心喜乐。这法律是以十诫赐给犹太人的，他们称之为十诫。这些诫命又总结为两条，即我们当全心、全灵、全意爱天主，并当爱人如己。因为全部法律和先知都系于这两条诫命，主自己已在福音中宣告，并以祂的榜样显示。同样，以色列子民从他们首次以某种形式庆祝逾越节，宰杀并吃那羔羊，用其血涂抹门框以确保安全的那一天起——从那天起，满了五十天，他们便领受了由天主的手指所写的法律，在这短语下我们已说过圣神被意指。同样，在主——真正的逾越节——受难并复活之后，圣神在第五十天以位格方式被派遣给门徒：但如今不再是刻在石版上（象征心的坚硬），而是当他们聚集在耶路撒冷同一处时，忽然从天上来了一阵响声，好像暴风刮来，又出现散开好像火舌的舌头，落在他们每人头上，他们就说起各种语言来，以致所有到他们那里的人都各自听见自己的方言（因为在那城中，犹太人惯于从各地聚集，他们散居各地，学会了各民族的多种语言）；此后，他们放胆宣讲基督，并因祂的名行了许多奇迹，以至于当伯多禄经过时，他的影子碰触了一个死人，那人就复活了。\par

42. 但犹太人看见因那被他们出于恶意或无知而钉死的人的名行了如此大的奇迹，其中有的恼怒起来，迫害宣讲祂的宗徒们；另一些则相反，因这事实——即他们曾嘲笑为被他们制伏并征服的那一位，如今却因祂的名行如此奇事——而更加惊奇，于是悔改并归信，以致成千的犹太人信了祂。这些人不再向天主渴求暂时的利益和地上的国度，也不再以血肉的方式期盼应许的君王基督；而是以不朽的方式领悟并爱慕那一位——祂曾以可朽之身替他们忍受了他们加给祂的沉重苦难，并将蒙赦免所有罪的恩赐赐给他们，甚至包括流祂血的罪孽，又藉祂复活的榜样，揭示了不朽作为他们应从祂那里盼望并渴求的对象。因此，他们如今治死了旧人的地上贪欲，被灵性生命的新经验所燃烧，正如主在福音中所吩咐的，他们变卖一切所有，把价银放在宗徒脚前，好使宗徒按各人的需要分给每人；他们以基督徒的爱德彼此和谐共处，不声称任何事物为私有，而是一切公用，并一心一德向着天主。后来，这些人自身也在肉身方面遭遇到同属血肉的同胞犹太人的迫害，并被分散到各处，以致因他们的分散，基督得以更广泛地被宣讲，同时他们自己也成为他们主的忍耐的追随者。因为那以温良忍耐了他们的主，吩咐他们为祂的缘故以温良忍耐一切。\par

43. 在这些圣人的迫害者中，宗徒保禄也曾列于其中，他异常凶猛地迫害基督徒。但后来，他成了信徒和宗徒，并被派遣向外邦人宣讲福音，为基督的名忍受了比他从前反对基督的名时所行的更痛苦的事。此外，在建立教会于他所播下福音种子的万国中，他常郑重吩咐：既然这些归信者（因他们来自偶像崇拜，且不熟悉唯一天主的崇拜）不易以变卖和分施财产的方式事奉天主，他们应为在犹太地各教会中信了基督的圣人中的贫穷弟兄捐献。这样，宗徒的教导使一些人成为士兵，另一些人成为行省的纳贡者，同时将基督置于中心，如同基石（正如先知早已预言的），使双方——即来自犹太人和来自外邦人——像从不同方向推进的墙壁，在亲属般的情感中联结在一起。但后来，来自不信的外邦人对基督教会更重更频繁的迫害兴起，主所说的话日复一日应验：\BibleQuote{看，我派遣你们好像羊进入狼群中。}\par

% structure: p0072 source=h2 target=section reason=relative-heading-rank; base=h2

\section{第24章——论教会如同发芽并受修剪的葡萄树}

44. 但那葡萄树的枝子，正如预言的，也如主亲自预言的，遍满全地，结出累累果实；它因着殉道者更丰富的血潮浇灌，反而长得愈发茂盛。当这些人为信仰的真谛在无数地域丧生时，连那逼迫的国度本身也歇手了，在骄傲的颈项折断后，归向对基督的认识和敬拜。再者，按主多次的预告，这同一棵葡萄树应当被修剪，不结果的枝条要从它上面剪除，这些枝条在基督名下，在各处引发了异端和分裂，分裂出自那些寻求自己荣耀而非祂荣耀的人；然而，他们的对抗也越发锻炼了教会，考验并显明了她的教义和忍耐。\par

45. 这一切，如今我们正看见它们完全应验，正如我们在许久以前所发的预言中读到的那样。当初的基督徒，因未亲眼看见这些事在他们的时代实际成就，被奇迹所动而相信；至于我们，既然这一切如今都已照着那些书中所记载的——那些书写在这一切事成就之前很久，其中所有事项都被预告为未来之事，正如它们如今已被视为现实——应验了，我们便被建立于信德，因此我们忍耐并持守于主内，毫不迟疑地相信那些尚待成就的事也必按定旨完成。事实上，在同一圣经中，仍能读到将来的苦难，以及最后的审判之日；那时，这两个城邦的所有居民都将重获他们的身体，复活，并在基督的审判座前为自己的生活交账。因为祂将要带着祂权能的荣耀而来，祂昔日曾屈尊就卑，以人性降世；祂将要把所有虔敬的人从不虔敬的人中分别出来——不仅是从那些完全拒绝相信祂的人，也是从那些相信祂却无果效、不结果实的人。祂要给一类人赐予永恒国度与祂同享，而给另一类人赐予永恒惩罚与魔鬼同受。然而，正如任何暂时的快乐都无法与圣徒们注定要获得的永生之乐相提并论，同样，任何暂时的痛苦也无法与不义者永久的痛苦相比。\par

% structure: p0075 source=h2 target=section reason=relative-heading-rank; base=h2

\section{第25章——论在复活信仰中的恒心}

46. 所以，弟兄，你要在你所相信的那位的名字和帮助下坚固自己，以抵抗那些嘲笑我们信仰之人的舌头；在他们身上，魔鬼说出诱惑的话，尤其致力于嘲弄复活的信仰。但你要从你自己的历史来推断：既然你曾存在，将来也会存在，正如你如今看见自己存在一样，虽然从前你并不存在。因为在几年之前，在你出生之前，甚至在你于母腹中成胎之前，你这伟大的身体结构，以及这四肢的形成和连接，曾在哪里？我再说，那时你这身体的结构和身量在哪里？它岂非从创造的隐秘之处，在看不见的主天主的塑造运作下，来到光明之中，并随着生命相继的阶段，按照那些固定的增长尺度，达到现今的规模和形态吗？那么，对于天主来说，祂也能在一瞬间从隐秘处召聚云团，顷刻间遮蔽天空，使你的身体恢复到原来的状态，这难道有任何困难吗？既然祂能够使你的身体成为它从前所不是的样子。因此，你要以刚毅不动摇的精神相信：那些似乎从人的眼前消失而归于朽坏的一切，在天主的全能面前都是安全且免于损失的；祂要在祂乐意的时候，毫无迟延或困难地予以恢复——至少，我应该说，那些祂的公义判断为配得恢复的部分；为使人能在这事上凭他们的身体报告他们的行为；并且在这身体中，他们或配得按他们虔诚的功绩接受属天不朽的交换，或按他们邪恶的功绩接受身体的可朽状况——而且那种状况不是可以被死亡解除的，而是给永久的痛苦提供材料。\par

所以，你要以坚定的信德和善行逃避——逃避吧，兄弟，那些折磨，其中折磨者不匮乏，受折磨者也不死；对他们而言，在痛苦中不能死就是无终止的死亡。你当燃起对圣人们永生的爱与渴望，在那里，行动不会劳苦，安息不会懈怠；在那里，赞美天主不会令人厌倦，也没有缺陷；在那里，心灵没有疲惫，身体没有衰竭；在那里，你也不会缺乏，无论是你自己的缺乏而需救助，还是邻人的缺乏而急于救助。天主将是那座圣城的全部享受和满足，那城生活于祂并藉祂而生活，在智慧与福乐中。因为正如我们所盼望和期待的，藉祂所应许的，我们将与天主的众天使相等，并同他们一起，现在藉信德行走，那时我们将面对面地享见那天主圣三。因为我们相信我们未见之事，好藉这信德的功绩，我们得以配得看见我们所相信的，并常住其中；以致于圣父、圣子及圣神平等的奥秘，以及同一天主圣三的合一，和这三个位格如何是一位天主，不再需要我们用信德的言词和有声的音节来述说，而是在那寂静中，以最纯洁、最炽热的默观来啜饮。\par

你要把这些铭刻在心，并呼求你相信的天主，求祂保护你免于魔鬼的诱惑；要当心，免得那仇敌从意想不到的方向偷袭你——他因自己受罚而极其恶毒地寻求陪伴，要找他人一同受罚。他不仅敢于藉那些仇恨基督徒名声的人，或那些因看到世界被这名号占据而痛苦、仍想侍奉偶像和邪灵怪异仪式的人，来试探基督徒；有时他也藉我们刚才提到的那类人，即那些如同葡萄树修剪时被剪下的枝条、与教会分离的人，就是所谓异端者或裂教者，来进行试探。然而，有时他也藉犹太人做同样的尝试，试图藉他们来试探和诱骗信徒。然而，最需要防范的是，任何人都不应被那些在天主教会内部、像糠秕一样被托住直到扬场之时的人所试探和欺骗。因为天主容忍这些人，目的在于藉他们的乖戾来操练和坚固祂选民的信德和明智；同时祂也顾念到，他们中有许多人进步了，并带着强烈的动力悔改而承行天主的旨意，因为他们被引导怜悯自己的灵魂。因为并非所有的人都因天主的容忍而为自己积蓄忿怒，直到祂公义审判的忿怒之日；而是有许多人因全能者的容忍而被引导至最有益的痛悔。直到此事成就，他们不仅成为那些已持守正道者的忍耐的操练，也成为他们怜悯的工具。因此，你将看到许多酗酒者、贪婪者、欺骗者、赌徒、奸淫者、行邪淫者、佩带亵渎护身符者，以及沉迷于巫师、占星家和精通各种不敬虔法术的占卜者的人。你也会看到，那些在异教节日充满剧场的人群，同样在基督徒节日充满教堂。当你看到这些时，你就会被引诱去效仿他们。不，我何必说\Emph{你要看见}呢？因为你已经知道这些事。你并非不知道，许多被称为基督徒的人，也行我所提及的所有这些恶事。你也不知道，有时也许你所认识的那些称为基督徒的人，甚至犯下比这些更严重的罪行。但如果你以为处在安全的位置上，就可以行这些事，那你就大错了；当基督开始以最严格的审判时，基督这名号对你也毫无用处，祂昔日也曾以最大的慈悲屈尊来救助人。因为祂曾预言这些事，并在福音中这样说：\BibleQuote{不是凡向我说\Quote{主啊！主啊！}的人，就能进入天国；而是那承行我在天之父旨意的人。在那一天，有许多人要向我说：\Quote{主啊！主啊！我们不是因你的名字吃过喝过吗？}}因此，对一切坚持这类行为的人，结局就是灭亡。所以，当你看到许多人不仅行这些事，还辩护和推荐它们时，你要坚定遵守天主的法律，不要追随那些故意违犯者。因为审判你的，不是他们的心意，而是祂的真理。\par

\Quote{当与善人结交，因你见他们与你在爱慕你的君王上同心。因为你若开始培养那样的品性，便会发现有许多这样的善人。因为在公开的竞技中，你尚且希望与性情相投的人为伴，并紧密依附那些与你在喜爱某位车夫、某位猎人或某位演员上志同道合的人，那么，你更应当乐于与那些在爱慕天主上与你同心的人为伴；凡爱慕祂的人，永不会有羞愧脸红之时，因为祂不仅自身不可能被战胜，也要使那些满怀爱慕地对待祂的人成为不可战胜的。同时，即便对于那些先你或与你同行走向天主的善人，你也不应将希望寄托于他们，正如你也不应将希望寄托于你自己，无论你已取得多么大的进步，而应寄托于那使他们和你成义、并使你成为你所是的那位。因为你在天主内是安稳的，因祂不变；但在人身上，没有人会明智地自认为安稳。但若我们应当爱那些尚未成义的人，以使他们成为义人，那么对于那些已经成义的人，应当爱得多么热烈呢？同时，爱一个人与将希望寄托于一个人是两回事；这区别如此之大，以至于天主命令前者而禁止后者。此外，若你因基督之名的缘故而遭受任何侮辱或苦难，既不从信仰上跌倒，也不从善道上偏离，你必获得更大的赏报；而那些在这样的情形下向魔鬼屈服的人，连较小的赏报也失掉了。但在天主面前要谦卑，以免祂允许你受诱惑超过你的力量。}\par

% structure: p0080 source=h2 target=section reason=relative-heading-rank; base=h2

\section{第二十六章——论慕道者的正式录取，及其所用标记}

50. 在结束这番话之后，应询问那人是否相信这些事并切愿遵守它们。当他如此回答后，当然要按教会的惯例庄严地为他划十字并予以处理。关于他所领受的圣事，首先应让他深切注意：神圣之事的标记固然是可见之物，但不可见之事本身也在其中受到尊崇；而那被祝福所圣化的形质，因此不应被视为如同在普通用途中看待的那样。此后，应告诉他他所聆听的言语形式所象征的是什么，以及那物质形态所代表的恩宠在他内所调和的又是什么。接着，我们应借那仪式劝告他：若他在圣经中听到任何带有属世含义的言辞，即使他不理解，也应相信其中必有所象征，指向圣德的品性和将来的生命。此外，他应简要地学到：凡他在正典书籍中听到的任何无法归向对永恒、真理、圣德的爱以及对近人的爱的话语或事迹，他都应相信是以比喻的方式说或做的；因此，他应努力理解，使之归向那双重爱德。但他还应进一步受劝告：不要按属世的意义理解\Emph{邻人}一词，而应理解为凡将来在圣城中与他同在的人，无论已在那里或尚未显现。最后应劝告他：不要对任何人的悔改绝望，因为他看到那人活在天主的容忍之下，别无他因，正如宗徒所说，乃是为使他悔改。\par

51. 若这篇我假设自己当面教导某位未受训者的讲话，在你看来太长，你可以更简短地阐述这些事。然而，我不认为它应比这更长。同时，很大程度上取决于事情本身的发展所可能提供的建议，以及听众实际所表现出的不仅是忍受，还有渴望。但当需要快速处理时，请注意整个事情是多么容易解释。假设又有一位愿意成为基督徒的人来到我们面前；并进一步假设，我们已询问过他，而他已做了我们假定那位先前的慕道者所做的回答；即使他拒绝这样回答，至少必须说他本应如此回答——那么，剩下要对他说的所有话，应以如下方式归纳：\par

52. 弟兄，确实，那应许给圣徒的将来世界中的福乐是伟大而真实的。反之，一切可见之物都要过去，这世界的一切虚荣、享乐和忧虑都将消亡，且它们（现在）将那些爱慕它们的人一同拖向毁灭。慈悲的天主愿意拯救人脱离这毁灭，即脱离永罚——只要他们不与自己为敌，不抗拒造物主的慈悲——就派遣了祂的独生\OriginalTerm{圣子}{Son}，即祂的\OriginalTerm{圣言}{Word}，与祂同等，祂藉着祂创造了万有。祂既住在祂的天主性内，既不离开圣父，也毫无改变，同时却取了人性，以可朽的肉身向人显现，来到人中间；好使正如死亡藉着一个人，即受造的第一个人亚当，进入人类——因为他同意妻子被魔鬼诱惑以致他们违犯了天主的诫命——同样，也藉着一个人，即耶稣基督，祂也是天主，天主子，所有相信祂的人过去的一切罪过都得以消除，并进入永生。\par

% structure: p0084 source=h2 target=section reason=relative-heading-rank; base=h2

\section{第二十七章 论旧约预言在教会中的可见应验}

53. 因为所有你现在在天主的教会中所见证的、并在基督的名下于全世界发生的事，早在许多世代前就已预言。我们读到了这些预言，如今也亲眼看见。藉着这些事，我们的信德得以建立。古时曾有一次淹没全地的洪水，其目的是为毁灭罪人。然而，凡在方舟里得救的人，都预示了将来的教会，这教会如今在世界之波上漂浮，并藉基督十字架的木获得拯救，不被淹没。曾向天主的忠仆\Person{亚巴郎}，一个独身的人，预言从他那里要生出一个民族，在列邦中，他们应在一切拜偶像的民族中敬拜独一真主；所有关于那民族命运所预言的，都照所预言的成就了。在那民族中，基督，诸圣的君王和他们的天主，也曾被预言要按肉身从同一\Person{亚巴郎}的后裔而来，那肉身祂取为己有，为使所有追随祂信德的人也成了亚巴郎的子孙；这事就这样成就了：基督由\Person{童贞玛利亚}诞生，她属于那民族。先知们预言祂将因那同一犹太民族的手而受苦于十字架上，而祂按肉身也正是从那民族而来；这事就这样成就了。预言祂将复活：祂已复活；并按照同一先知的预言，祂升了天，并派遣圣神给祂的门徒。不仅先知们，连主耶稣基督自己也预言了祂的教会将遍及全世界，因圣人的殉道和受苦而扩展；这预言是在祂的名字尚未在外邦人中被宣告，且在被知晓之处受到嘲笑时说的；然而，藉着祂神迹的能力——无论是祂亲手所行的，还是藉祂仆人所行的——当这些事被传扬和相信时，我们已看到所预言的应验，并且看见地上的君王们，从前惯于迫害基督徒，如今也屈服于基督的名下。也曾预言从祂的教会中会产生分裂和异端，并假借祂的名在他们所能支配之处寻求自己的荣耀而非基督的荣耀；这些预言都已应验。\par

54. 那么，那些尚未成就的事岂会不实现呢？显然，正如先前预言的那些事已经实现，后者也必将实现。我指的是所有义人尚待应验的磨难，以及审判日，那日将在死人复活中把一切恶人与义人分开——不仅将教会之外的恶人如此分开，而且将教会本身的糠秕分离出来，投入应受的火中，这糠秕必须以极大的忍耐容忍到最后簸扬。此外，那些嘲笑（死人）复活的人，因为他们认为这肉体既会朽坏，就不能复活，他们必将在同一肉体中复活而受惩罚，天主将使这样的人明白：那在未有形体时能形成这些身体者，也能在顷刻之间恢复它们原来的样子。但所有注定与基督一同作王的信徒，将带着同一身体复活，并且被恩赐改变为天使般的不朽，使他们能与天主的众天使相等，正如主自己所应许的；并且要无间断、无疲惫地赞美祂，永远生活在祂内并藉祂而生活，充满人既无法表达也无法理解的喜乐和真福。\par

\Quote{因此，你要相信这些事，并谨防诱惑（因为魔鬼寻求他人与他一同丧亡）；这样，那仇敌不仅无法藉着教会之外的人——无论是异教徒、犹太人或异端者——来诱惑你，而且你本人也不愿效法天主教会内部那些你看见的度邪恶生活的人，他们或在肚腹和喉咙的享乐上放纵无度，或不贞洁，或沉溺于好奇迷信的虚妄与非法仪式，无论他们是沉湎于公共表演（的空虚）、邪术或魔鬼的占卜，还是生活在贪婪与骄傲的浮华和膨胀中，或追求任何法律所谴责和惩罚的生活方式。反之，你要与善人联合；你若自己也成为善人，便容易找到他们，好使你们彼此结合，为天主本身而崇拜祂、爱慕祂，因为祂自身将是我们的圆满赏报，为使我们能在那生命中享受祂的良善与美善。然而，爱祂的方式不应像爱任何用眼睛可见的对象那样，而应如爱智慧、真理、圣洁、正义、爱德，以及任何类似性质的事物；并且，这种爱不是指向这些事物在人与上的表现，而是指向它们在不朽不变智慧的泉源本身。因此，你若看见何人爱慕这些事物，就要依附他们，好藉着那降生成人、成为天主与人之间中保的基督，与天主和好。至于那些心术不正的人，即使他们进入教会的围墙内，也不要以为他们能进入天国；因为到了时候，他们若不改恶向善，就会被分别出来。所以，你要效法善人的榜样，容忍恶人，爱所有的人；因为你不知道今天为恶的人明天会成为怎样的人。然而，你不要爱他们的不义，却要爱他们本人，明确目的在于使他们能领会正义；因为不仅爱天主的命令加于我们，爱邻人的命令也是如此，这两条诫命是全部法律和先知的基础。除非一个人领受了另一项恩赐——圣神（祂与圣父、圣子同等），没有人能满全这一切；因为这位天主圣三就是天主，一切希望应寄託于这位天主。人的希望不应寄託于任何人，不论他是什么样的人。因为使我们成义的那一位是一回事，与我们一起受成义的众人是另一回事。此外，魔鬼不仅藉着情欲来诱惑，也藉着侮辱、痛苦甚至死亡本身的恐怖来诱惑。但任何人若因基督的名和永生的希望而受苦，并以恒心忍受，他将得到更大的赏报；反之，若他向魔鬼屈服，他将与魔鬼一同受罚。然而，慈悲的工作与虔诚的谦卑相结合，会从天主那里得到这样的承认：祂必不让祂的仆人受诱惑超过他们所能承受的。}\par

\SourceRecord{Translated by S.D.F. Salmond.；Nicene and Post-Nicene Fathers, First Series；Vol. 3.；Edited by Philip Schaff.；Buffalo, NY: Christian Literature Publishing Co.,；1887.；<http://www.newadvent.org/fathers/1303.htm>.}
